\documentclass[reqno]{amsart}
\usepackage{a4wide}
\usepackage{hyperref}
\usepackage{mathrsfs}

\AtBeginDocument{{\noindent\small
\emph{Electronic Journal of Differential Equations},
Vol. 2026 (2026), No. 05, pp. 1--18.\newline
ISSN: 1072-6691. URL: https://ejde.math.txstate.edu, DOI: 10.58997/ejde.2026.05}
\thanks{\copyright 2026. This work is licensed under a CC BY 4.0 license.}
\vspace{8mm}}

\begin{document}
\title[\hfilneg EJDE-2026/05\hfil Multiple solutions for Kirchhoff type systems]
{Multiple solutions for Kirchhoff type systems involving singular and critical nonlinearities}

\author[M. Louchaich \hfil EJDE-2026/05\hfilneg]
{Mohamed Louchaich}

\address{Mohamed Louchaich \newline
High school of sciences and technology, 
Sousse University,
4011 Hammam Sousse, Tunisia}
\email{mohamedlouchaiech@gmail.com}

\thanks{Submitted September 20, 2025. Published January 17, 2026.}
\subjclass[2020]{35A15, 35D30,  35J60, 35J75, 35R11, 46E35, 47G20}
\keywords{p-Laplacian; Kirchhoff system; critical point theory; multiple positive solution;
\hfill\break\indent singular nonlinearity, critical Sobolev exponent}

\begin{abstract}
This article investigates the fractional singular Kirchhoff system
\begin{gather*}
m(\mathcal{N}(\phi,\psi))\mathcal{L}_\mathcal{K}^p(\phi)
=\lambda a(x)\phi^{-\gamma_1}+\frac{\theta_1}{p^\star_{N,s}}\phi^{\theta_1-1}
\psi^{\theta_2} \quad \text{in } \mathcal{D} \\
m(\mathcal{N}(\phi,\psi))\mathcal{L}_\mathcal{K}^p(\psi)
=\lambda b(x)\psi^{-\gamma_2}+\frac{\theta_2}{p^\star_{N,s}}\phi^{\theta_1}
 \psi^{\theta_2-1} \quad \text{in } \mathcal{D} \\
\phi>0,\quad \psi>0  \quad\text{in } \mathcal{D} \\
\phi=\psi=0 \quad \text{in } \mathbb{R}^N\setminus \mathcal{D},
\end{gather*}
where
$$
\mathcal{N}(\phi,\psi)=\int_{\mathbb{R}^{2N}}|\phi(x)-\phi(y)|^p \mathcal{K}(x,y)\,dx\,dy
+\int_{\mathbb{R}^{2N}}|\psi(x)-\psi(y)|^p\mathcal{K}(x,y)\,dx\,dy.
$$
Here, $\mathcal{D}$ is a bounded domain in $\mathbb{R}^N$ with a
Lipschitz boundary  $\partial\mathcal{D}$.
$\mathcal{L}_\mathcal{K}^p$ is a non-local operator with
a singular kernel $\mathcal{K}$.
$p> 1$, $\lambda>0$, and $m$ is a continuous function.
$\gamma_1,\gamma_2\in (0,1)$. and $a,b$ are  non-negative bounded functions.
$\theta_1,\theta_2>1$ and $\theta_1+\theta_2=p^\star_{N_s}$,
where $p^\star_{N_s}$ is the fractional critical Sobolev exponent
$p^\star_{N_s}=\frac{Np}{N-sp}$.
Our findings encompass the degenerate case in the fractional
setting, allowing the Kirchhoff function $m$ to take zero value at zero.
We employ Kajikiya's version of the symmetric mountain pass lemma to prove the
existence of a sequence of infinitely many small solutions with negative energy
that converge to zero.
\end{abstract}

\maketitle
\numberwithin{equation}{section}
\newtheorem{theorem}{Theorem}[section]
\newtheorem{lemma}[theorem]{Lemma}
\newtheorem{definition}[theorem]{Definition}
\newtheorem{proposition}[theorem]{Proposition}
\allowdisplaybreaks

\section{Introduction and main result}

Recently, the study of non-local equations and systems has attracted a lot of 
attention. Various studies have been considered on this topic, see 
\cite{t1, t2, t3, t4, D, t5, a, t7, s, qui, serv2, serv4, t9, t10}. 
Nonlocal equations have been widely used in many fields of sciences, 
including continuum mechanics, phase transition phenomena, population dynamics, 
and game theory, especially those involving fractional and nonlocal elliptic operators. 
As discussed in \cite{l1, caff}, these operators naturally appear as stochastic 
stabilizers of L\'evy processes..

On the other hand, a great deal of attention has been focused on nonlocal fractional 
equations with critical nonlinearities. This area of study has received significant 
attention in recent years leading to several advancements and investigations.
Among the references we like to mention \cite{Alv, ambr,t1,daoua,f2,o,Mok, o2,p,t9}. 
These papers provide valuable insights and analysis of nonlocal fractional equations 
with critical nonlinearities, and they serve as important references for further 
exploration in this area.

In this work, we study the existence of solutions for a fractional Kirchhoff system 
that involve singular and critical nonlinearities. 
The system studied is
\begin{equation}\label{System}
\begin{gathered}
m(\mathcal{N}(\phi,\psi))\mathcal{L}_\mathcal{K}^p(\phi)
=\lambda a(x)\phi^{-\gamma_1}+\frac{\theta_1}{p^\star_{N,s}}\phi^{\theta_1-1}
\psi^{\theta_2} \quad \text{in } \mathcal{D} \\
m(\mathcal{N}(\phi,\psi))\mathcal{L}_\mathcal{K}^p(\psi)
=\lambda b(x)\psi^{-\gamma_2}+\frac{\theta_2}{p^\star_{N,s}}\phi^{\theta_1}
\psi^{\theta_2-1} \quad \text{in } \mathcal{D} \\
\phi>0,\quad \psi>0  \quad\text{in } \mathcal{D} \\
\phi=\psi=0 \quad \text{in } \mathbb{R}^N\setminus \mathcal{D},
\end{gathered}
\end{equation}
where
$$
\mathcal{N}(\phi,\psi)=\int_{\mathbb{R}^{2N}}|\phi(x)-\phi(y)|^p \mathcal{K}(x,y)\,dx\,dy
+\int_{\mathbb{R}^{2N}}|\psi(x)-\psi(y)|^p\mathcal{K}(x,y)\,dx\,dy,
$$
$\mathcal{D}$ is a smooth bounded domain in $\mathbb{R}^N$, $\lambda>0$ is a real 
parameter, $p>1$, $\gamma_1,\gamma_2\in (0,1)$, and $\theta_1,\theta_2>1$ satisfy 
$\theta_1+\theta_2=p^\star_{N,s}$. 
The fractional critical Sobolev exponent $p^\star_{N,s}$ is defined as 
$p^\star_{N,s}=\frac{Np}{N-sp}$. 
$a,b\in L^\infty(\mathcal{D})$, $a>0$, $b>0$ a.e. in $\mathcal{D}$. 
$\mathcal{L}_\mathcal{K}^p$ is a non-local operator defined for any smooth functions 
$\phi:\mathbb{R}^N\to\mathbb{R}$ by
$$
\mathcal{L}_\mathcal{K}^p\phi(x)
=2\lim_{\epsilon\to 0}\int_{\mathbb{R}^N\setminus B_\epsilon(x)}
|\phi(x)-\phi(y)|^{p-2}(\phi(x)-\phi(y))\mathcal{K}(x,y)\,dy,\quad x\in \mathbb{R}^N.
$$
Here, $\mathcal{K}:\mathbb{R}^{2N}\setminus \{0,0\}\to (0,+\infty)$ is a
 measurable function satisfying the following properties:
\begin{gather*}
\mathcal{K}(x,y)\geq \mathcal{K}_0|x-y|^{-(N+sp)}\quad
 \text{for all } (x,y)\in \mathbb{R}^{2N},\; x\neq y\\
\sigma\mathcal{K}\in L^1(\mathbb{R}^{2N})\quad \text{with }
 \sigma(x,y)=\min\{1,|x-y|^p\}\\
\mathcal{K}(x,y)=\mathcal{K}(y,x)\quad\text{for all } (x,y)\in \mathbb{R}^{2N}.
\end{gather*}
A commonly used model for the kernel function $\mathcal{K}$ is given by the singular 
kernel $\mathcal{K}(x,y)=|x-y|^{-(sp+N)}$. In this case, the operator 
$\mathcal{L}_\mathcal{K}^p$ corresponds to the fractional $p$-Laplacian, which can be 
defined (up to normalization factors) for any $\phi\in C^\infty_0(\mathbb{R}^N)$ by:
$$
(-\Delta_p)^s\phi(x)
=2\lim_{\epsilon\to 0}\int_{\mathbb{R}^N\setminus B_\epsilon(x)}
\frac{|\phi(x)-\phi(y)|^{p-2}(\phi(x)-\phi(y))}{|x-y|^{sp+N}}\,dy,\quad x\in\mathbb{R}^N.
$$
The Kirchhoff function $m$ is assumed to satisfies the following assumptions:
\begin{itemize}
 \item[(A1)] For each $\delta>0$, there exists $\kappa_\delta>0$ such that
 $m(z)\geq \kappa_\delta$ for all $z\geq \delta$.

\item[(A2)] There exists $\sigma\in [1,\frac{p^\star_{N,s}}{p})$ such that
 $$ 
\sigma\mathscr{M}(z)=\sigma\int_0^z m(\varsigma)d\varsigma
\geq z m(z),\quad \forall \xi\in \mathbb{R}_+.
$$

\item[(A3)] There exists $m_0>0$ such that
$m(z)\geq m_0z^{\sigma-1}$ for all $\xi\in [0,1]$.
\end{itemize} 
A prototype for $m$, proposed by Kirchhoff, is 
\begin{equation}\label{Kirchhoff}
m(z)=\alpha+\beta z^{\sigma-1},\quad \alpha,\beta \geq 0,\quad \alpha+\beta>0,
\quad \sigma>1.
\end{equation}

The Kirchhoff problem \eqref{System} is classified as non-degenerate if 
$m(z)\geq \overline{m}>0$ for all $z\in \mathbb{R}_+$.
A non-degenerate case can be achieved, for example, when $\alpha>0$ and $\beta\geq 0$
in the model case \eqref{Kirchhoff}. For recent results on non-degenerate Kirchhoff-
type problems, we refer the reader to \cite{o, Ny, Ny1, Ny2, p}. On the other hand, 
if $m(0)=0$ but $m(z)>0$ for all $z>0$, the Kirchhoff problem is considered degenerate. 
This degenerate scenario occurs in the model case  \eqref{Kirchhoff} when 
$\alpha=0$ and $\beta>0$. Relevant research papers on degenerate Kirchhoff-type problems
include \cite{Au, Col, daoua, Liu, s}.

Fiscella and Valdinoci  \cite{fisc} presented a detailed examination of the physical 
understanding of fractional Kirchhoff problems and their applications.  
They considered a stationary Kirchhoff  problem that capture the non-local component 
of tension caused by nonlocal measurements of a string's fractional length.  
In this case, the function $m$ measures the change in tension on the string produced 
by differences in its length during vibration.  
The fact that $m(0)=0$ implies that the string's base tension is zero, indicating 
a realistic model.

The aim of this article is to investigate the existence of solutions for 
system \eqref{System} by  using variational methods and critical point theory. 
Before presenting our main result, we introduce the following notation and
functional setting. Let us define the space
$$
E=\Big\lbrace \phi\in L^p(\mathbb{R}^N) \text{ and }
  \int_{\mathbb{R}^{2N}}|\phi(x)-\phi(y)|^p \mathcal{K}(x,y)\,dx\,dy
  <\infty\Big\rbrace,
$$
equipped with the norm
$$
\| \phi\|_{E}=\Big(\| \phi\|_{L^p(\mathbb{R}^N)}^{p}
+\int_{\mathbb{R}^{2N}}|\phi(x)-\phi(y)|^p \mathcal{K}(x,y)\,dx\,dy
\Big)^{1/p}.
$$
We also introduce the  space
$$
E_0(\mathcal{D})=\Big\lbrace  \phi\in E:
 \phi=0 \text{ a.e. }  \mathbb{R}^N\setminus\mathcal{D}\Big\rbrace,
$$
where $\mathcal{D}$ is a given domain.
 According to \cite[Lemma 4]{z}, $E_0(\mathcal{D})$ is a separable and reflexive 
 Banach space, which can be equipped with the  norm
\begin{equation*}
\| u\|_{E_0}=\Big(\int_{\mathbb{R}^{2N}}|\phi(x)-\phi(y)
|^p \mathcal{K}(x,y)\,dx\,dy\Big)^{1/p},\quad \forall \phi\in E_0(\mathcal{D}).
\end{equation*}

Given that $\mathcal{D}$ is a smooth bounded domain, it is well-known that the embedding $E_0(\mathcal{D})\hookrightarrow L^\mu(\mathcal{D})$ holds continuously for $\mu\in [1,p^\star_{N,s}]$ (refer to \cite{nez}, Lemma 2.3) and compactly for $\mu\in [1,p^\star_{N,s})$, where $p^\star_{N,s}=\frac{Np}{N-sp}$. Additionally, there exists a positive constant $C_\mu$ such that for any $\phi\in E_0(\mathcal{D})$, the following inequality holds:
\begin{equation}\label{sobolevinequaity}
\| \phi\|_{L^\mu(\Omega)}\leq C_\mu \| \phi\|_{E_0}.
\end{equation}
In what follows, $S_\star$ stands for the optimal constant for the Sobolev embedding 
$E_0(\mathcal{D})\hookrightarrow L^{p^\star_{N,s}}(\mathcal{D})$, 
which is given by
\begin{equation}\label{S1}
S_\star=\inf_{\phi\in E_0(\mathcal{D})\setminus \lbrace 0\rbrace}
\frac{{\big(\int_{\mathbb{R}^{2N}}|\phi(x)-\phi(y)|^p\mathcal{K}(x,y)\,dx\,dy
\big)^{1/p}}}{{\big(\int_\mathcal{D} |\phi(x)|^{p^\star_{N,s}}\,dx
\big)^{1/p^\star_{N,s}}}}.
\end{equation}
Let $X=E_0(\mathcal{D})\times E_0(\mathcal{D})$, equipped with the usual norm 
$$
\| (\phi,\psi)\|=\Big(\| \phi\|_{E_0}^p+\| \psi\|_{E_0}^p\Big)^{1/p},
\quad \text{for } (\phi,\psi)\in X.
$$

\begin{definition}\rm
We say that $(\phi,\psi)\in X$ is a weak solution of system \eqref{System} if
\begin{equation} \label{p11}
\begin{aligned}
&m(\mathcal{N}(\phi,\psi))\int_{\mathbb{R}^{2N}}|\phi(x)
 -\phi(y)|^{p-2}(\phi(x)-\phi(y))(\varphi_1(x)
 -\varphi_1(y))\mathcal{K}(x,y)\,dx\,dy \\
&+m(\mathcal{N}(\phi,\psi))\int_{\mathbb{R}^{2N}}|\psi(x)
-\psi(y)|^{p-2}(\psi(x)-\psi(y))(\varphi_2(x)-\varphi_2(y))\mathcal{K}(x,y)\,dx\,dy \\
&=\lambda \int_\mathcal{D} a(x)\phi^{-\gamma_1}\varphi_1\,dx+ \int_\mathcal{D}b(x)
 \psi^{-\gamma_2}\varphi_2\,dx \\
&+\frac{\theta_1}{p^\star_{N,s}}\int_\mathcal{D}\phi^{\theta_1-1}
\psi^{\theta_2}\varphi_1\,dx
+\frac{\theta_2}{p^\star_{N,s}}\int_\mathcal{D}\phi^{\theta_1}\psi^{\theta_2-1}
\varphi_2\,dx,\quad \forall (\varphi_1,\varphi_2)\in X.
\end{aligned}
\end{equation}
\end{definition}

To obtain the existence of  weak solutions for system \eqref{System}, 
we seek critical points of the associated energy functional 
$\mathcal{J}_\lambda:X\to\mathbb{R}$ given by
$$
\mathcal{J}_\lambda(\phi,\psi)=\frac{1}{p}\mathscr{M}(\mathcal{N}(\phi,\psi))
-\lambda \Phi(\phi,\psi)-\Psi(\phi,\psi), 
$$
where
$$
\Phi(\phi,\psi)=\frac{1}{1-\gamma_1}\int_\mathcal{D} a(x)(\phi^+)^{1-\gamma_1}\,dx
+\frac{1}{1-\gamma_2}\int_\mathcal{D}b(x)(\psi^+)^{1-\gamma_2}\,dx, 
$$
and
$$
\Psi(\phi,\psi)=\frac{1}{p^\star_{N,s}}\int_\mathcal{D}  
(\phi^+)^{\theta_1}(\psi^+)^{\theta_2}\,dx.
$$
Our main result can be summarized as follows.

\begin{theorem}\label{Theorem} 
Under  assumptions {\rm (A1)--(A3)} there is a positive constant $\lambda_0>0$ such 
that for any $\lambda\in (0,\lambda_0)$, system \eqref{System} possesses a sequence
 of weak solutions $(\phi_k,\psi_k)_{k\in\mathbb{N}}\subset X$ such that
 $\mathcal{J}_\lambda(\phi_k,\psi_k)< 0$ and $\|(\phi_k,\psi_k)\|\to 0$
 as $k\to +\infty$.
\end{theorem}

The remainder of this paper is organized as follows:
 In section 2, we present some preliminary results. 
 In section 3, we show the existence of infinitely many small weak solutions for 
 an associated approximating problem. 
 In section 4, we prove our main result.
 
\section{Preliminaries}

In this section, we give some important results that will be useful in the proof of 
our main result.

\begin{lemma} 
For each $(\phi,\psi)\in X$, there exists $S>0$ such that
\begin{equation}\label{S}
\| (\phi,\psi)\|\geq S \int_\mathcal{D}
|\phi(x)|^{\theta_1}|\psi(x)|^{\theta_2}\,dx.
\end{equation}
\end{lemma}

\begin{proof}
Let $(\phi,\psi)\in X$. By using  H\"older's inequality, one has
\begin{align*}
\int_\mathcal{D} |\phi(x)|^{\theta_1}|\psi(x)|^{\theta_2}\,dx
&\leq  \Big(\int_\mathcal{D}|\phi(x)|^{p^\star_{N,s}}\,dx
 \Big)^\frac{\theta_1}{p^\star_{N,s}}\Big(\int_\mathcal{D}|\psi(x)|^{p^\star_{N,s}}\,dx\Big)^\frac{\theta_2}{p^\star_{N,s}} \\
&\leq S_\star^{-(\theta_1+\theta_2)}\| \phi\|_{E_0}^{\theta_1}
 \| \psi\|_{E_0}^{\theta_2}\\
&\leq  \frac{1}{S}\| (\phi,\psi)\|^{p^\star_{N,s}},
\end{align*}
where $S=S^{\theta_1+\theta_2}$.
\end{proof}

In the following Lemma, we establishes a key result related to the convergence of 
a sequence in $E_0(\mathcal{D})$ to a limit within the same space.

\begin{lemma}\label{l1}
Let $(\phi_n)_{n\in\mathbb{N}}\subset E_0(\mathcal{D})$ be such $\phi_n\to \phi$
weakly in $E_0(\mathcal{D})$ for some $\phi\in E_0(\mathcal{D})$. 
Then, we have
\[
\| \phi_n-\phi\|_{E_0}^p+\| \phi\|_{E_0}^p
=\| \phi_n\|_{E_0}^p+o_n(1).
\]
\end{lemma}

\begin{proof}
We use a result of Br\'ezis-Lieb's lemma \cite{bb}. 
It states that if we have a bounded sequence $(\Phi_n)_{n\in\mathbb{N}}$ in the space 
$L^p(\mathbb{R}^d)$, where $p>1$, and $\Phi_n$ converges to $\Phi$ almost everywhere 
in $\mathbb{R}^d$, then we have
\[
|\Phi_n-\Phi|_{L^p(\mathbb{R}^d)}^p+|\Phi|_{L^p(\mathbb{R}^d)}^p
=|\Phi_n|_{L^p(\mathbb{R}^d)}^p+o_n(1).
\]
We define the function 
\[
\Phi_n=(\phi_n(x)-\phi_n(y))\sqrt[p]{\mathcal{K}(x,y)}, \quad \text{where } d=2N.
\]
Using this definition, we can write 
\begin{align*}
&\int_{\mathbb{R}^{2N}}|(\phi_n-\phi)(x)-(\phi_n-\phi)(y)|^p\mathcal{K}(x,y)
\,dx\,dy+\int_{\mathbb{R}^{2N}}|\phi(x)-\phi(y)|^p\mathcal{K}(x,y)\,dx\,dy \\
&=\int_{\mathbb{R}^{2N}}|\phi_n(x)-\phi_n(y)|^p\mathcal{K}(x,y)\,dx\,dy+o_n(1).
\end{align*}
This completes the proof.
\end{proof}

\begin{lemma}\label{l2} 
Let $\alpha,\beta\in (1,+\infty)$. Given $\epsilon>0$, there exists $C>0$ such that 
for each $\phi_1,\phi_2,v_1,v_2$, we have
\[
\Big||\phi_1+\psi_1|^\alpha|\phi_2+\psi_2|^\beta
-|\phi_1|^\alpha|\phi_2|^\beta\Big|
\leq \epsilon |\phi_1|^\alpha|\phi_2|^\beta+C(|\phi_1|^\alpha|\psi_2|^\beta
+|\psi_1|^\alpha|\phi_2|^\beta+|\psi_1|^\alpha|\psi_2|^\beta).
\]
\end{lemma}

\begin{proof}
Fix $0<\eta<\min\lbrace 1,\frac{\epsilon}{1+2^\alpha}  \rbrace$. 
Then there exists $\overline{C}>0$ such that
\begin{align*}
&\Big||\phi_1+\psi_1|^\alpha|\phi_2+\psi_2|^\beta-|\phi_1|^\alpha|\phi_2|^\beta\Big|\\
&\leq|\phi_1+\psi_1|^\alpha \Big||\phi_2+\psi_2|^\beta-|\phi_2|^\beta\Big|
 +\Big||\phi_1+\psi_1|^\alpha-|\phi_1|^\alpha|\phi_2|^\beta\Big|\\
&\leq 2^\alpha(|\phi_1|^\alpha+|\psi_1|^\alpha)(\eta|\phi_2|^\beta
 +\overline{C}|\psi_2|^\beta)
 +(\eta|\phi_1|^\alpha+\overline{C}|\psi_1|^\alpha)|\phi_2|^\beta,\\
&\leq \epsilon |\phi_1|^\alpha|\phi_2|^\beta+2^\alpha\overline{C}|
 \phi_1|^\alpha|\psi_2|^\beta
 +(2^\alpha+\overline{C})\psi_1|^\alpha|\phi_2|^\beta+2^\alpha\overline{C}\psi_1|^\alpha|\psi_2|^\beta,
\end{align*}
as claimed.
\end{proof}

The following lemma presents a key result regarding the convergence properties of a 
sequence $(\phi_n,\psi_n)$ in the space $X$, under weak convergence and almost 
everywhere convergence in $\mathbb{R}^N$.

\begin{lemma}\label{l3} 
Let $(\phi_n,v_n)\rightharpoonup (\phi,v)$ weakly in $X$ and 
$(\phi_n,v_n)\to (\phi,v)$ a.e. in $\mathbb{R}^N$. Then, for each $\alpha,\beta>1$
with $\alpha+\beta\leq p^\star_{N,s}$, up to a subsequence, we have
\begin{align*}
&\int_{\mathbb{R}^N}|\phi_n(x)-\phi(x)|^\alpha|\psi_n(x)-\psi(x)|^\beta\,dx
 +\int_{\mathbb{R}^N} |\phi(x)|^\alpha|\psi(x)|^\beta\,dx \\
&=\int_{\mathbb{R}^N}|\phi_n(x)|^\alpha|\psi_n(x)|^\beta\,dx+o_n(1).
\end{align*}
\end{lemma}

\begin{proof}
Passing to a subsequence, we have that
\begin{gather*}
\phi_n\to \phi,\ \psi_n\to \psi,\quad \text{a.e. in } \Omega,\\
\phi_n\to \phi\quad \text{in}  L^\alpha_{\rm loc}(\mathbb{R}^N), \quad\text{and}\quad
\psi_n\to \psi\quad \text{in } L^\alpha_{\rm loc}(\mathbb{R}^N).
\end{gather*}
Let $\epsilon>0$, and let $C>0$ be a fixed real number as in Lemma \ref{l2}. We define
\begin{equation} \label{u0}
\begin{aligned}
W_n&=\Big||\phi_n|^\alpha|\psi_n|^\beta-|\phi_n-\phi|^\alpha|\psi_n-\psi|^\beta
 -|\phi|^\alpha|\psi|^\beta\Big|-\epsilon |\phi_n-\phi|^\alpha|\psi_n-\psi|^\beta  \\
&\quad -C(|\phi|^\alpha|\psi_n-\psi|^\beta+|\phi_n-\phi|^\alpha|\psi|^\beta).
\end{aligned}
\end{equation}
It is easy to see that, $w_n\to 0$ a.e. in $\mathbb{R}^N$ and $w_n\in L^1(\mathbb{R}^N)$.
Furthermore, utilizing Lemma \ref{l2} with $\phi_1=\phi_n-\phi$, $\phi_2=\psi_n-\psi$,
$\psi_1=\phi$, and $\psi_2=\psi$, we obtain
\[
W_n\leq (C+1) |\phi|^\alpha|\psi|^\beta.
\]
Then, we apply the dominated convergence theorem to get
\begin{equation}\label{u1}
\lim_{n\to +\infty}\int_{\mathbb{R}^N}W_n^+\,dx=0,
\end{equation}
where $W_n^+=\max\lbrace W_n,0\rbrace$. Fix $T>0$ large enough so that
\begin{equation}\label{u2}
C\int_{|x|\geq T}|\phi|^\alpha|\psi_n-\psi|^\beta dx\leq C |\psi_n-\psi|^\beta_{L^{p^\star_s}(\mathbb{R}^N)}\Big(\int_{|x|\geq T}|\phi|^{p^\star_s}\Big)^\frac{\alpha}{{p^\star_s}}<\epsilon,
\end{equation}
and
\begin{equation}\label{u3}
C\int_{|x|\geq T}|\phi_n-\phi|^\alpha|\psi|^\beta dx\leq C |\phi_n-\phi|^\alpha_{L^{p^\star_s}(\mathbb{R}^N)}\Big(\int_{|x|\geq T}|\psi|^{p^\star_s}\Big)^\frac{\beta}{{p^\star_s}}<\epsilon.
\end{equation}
Then, for $n$ large enough, we have
\begin{equation}\label{u4}
C\int_{|x|\leq T}|\phi|^\alpha|\psi_n-\psi|^\beta dx\leq C \max_{|x|\leq T}|\phi(x)|^\alpha \int_{|x|\leq T}|\psi_n-\psi|^\beta dx <\epsilon,
\end{equation}
and
\begin{equation}\label{u5}
C\int_{|x|\leq T}|\phi_n-\phi|^\alpha|v|^\beta dx\leq C \max_{|x|\leq T}|v(x)|^\beta \int_{|x|\leq T}|\phi_n-\phi|^\alpha dx<\epsilon.
\end{equation}
Now, combining equations \eqref{u0}-\eqref{u5}, we obtain
\begin{align*}
&\int_{\mathbb{R}^N}\Big||\phi_n|^\alpha|v_n|^\beta-|\phi_n-\phi|^\alpha|v_n-v|^\beta
 -|\phi|^\alpha|v|^\beta\Big|\,dx\\
&\leq \epsilon \int_{\mathbb{R}^N}|\phi_n-\phi|^\alpha|v_n-v|^\beta\,dx
+C\int_{\mathbb{R}^N}|\phi|^\alpha|v_n-v|^\beta+|\phi_n-\phi|^\alpha|v|^\beta\,dx
+\int_{\mathbb{R}^N} W_n^+\,dx,\\
&< \tilde{C}\epsilon,
\end{align*}
for $n$ large enough, as claimed.
\end{proof}

Now, we recall some of topological techniques introduced by Krasnoselskii in \cite{cc}. 
These methods have been widely recognized as powerful tools to obtain solutions to 
a variety of mathematical problems. Specifically, in our case, we aim to prove the 
existence of a sequence of small solutions. Let $(X,|\cdot|_X)$ be a real Banach space. 
We introduce the class $\Sigma$, which consists of all closed subsets 
$A\subset X\setminus{0}$ that exhibit symmetry with respect to the origin. 
In other words, if $u\in A$, then $-u\in A$. This symmetry property is a key 
characteristic of the sets in $\Sigma$.

Now,  to characterize the properties and complexity of sets in $\Sigma$, 
we introduce the notion of Krasnoselskii's genus.

\begin{definition}\label{genus} \rm
Let $A\in \Sigma$. The Krasnoselskii's genus $\gamma(A)$ of $A$ is defined as the 
smallest positive integer $n$ such that there exists an odd mapping 
$h\in C(A,\mathbb{R}^n)$ satisfying $h(x)\neq 0$ for all $x\in A$. 
If such an integer $n$ does not exist, we set $\gamma(A)=+\infty$. Furthermore, 
we define $\gamma(\emptyset)=0$.
\end{definition}

Next, we will outline the essential properties of the genus that will be 
employed throughout this work. For more comprehensive information on this topic, 
readers are encouraged to consult the reference \cite{cc}.

\begin{proposition}\label{prop}
Let $A$ and $B$ be symmetric closed subsets of $E$ that do not include the origin. 
The following properties hold
\begin{enumerate}
\item If there exists an odd continuous mapping from $A$ to $B$, then $\gamma(A)\leq \gamma(B)$.
\item If there is an odd homeomorphism from $A$ onto $B$, then $\gamma(A)= \gamma(B)$.
\item If $\gamma(A)<\infty$, then $\gamma(A\setminus B)=\gamma(A)-\gamma(B)$.
\item The $n$-dimensional sphere $\mathbb{S}^n$ has a genus of $(n+1)$ as a 
consequence of the Borsuk-Ulam Theorem.
\end{enumerate}
\end{proposition}

Now we present a variant of the symmetric mountain pass lemma, originally formulated 
by Kajikiya \cite{n}. This lemma offers a powerful tool for analysing critical points 
of functionals with symmetric features.

\begin{lemma}[Kajikiya's Variant of the Symmetric Mountain Pass Lemma] \label{kajikiya}
Let $X$ be an infinite-dimensional Banach space, and let 
$\mathcal{J}\in C^1(X,\mathbb{R})$ be a functional satisfying the following 
conditions
\begin{enumerate}
\item  $\mathcal{J}$ is even, bounded from below, and $\mathcal{J}(0)=0$.
\item  $\mathcal{J}$ satisfies the local Palais-Smale condition, which means 
 that for some $\overline{c}\in\mathbb{R}$, any sequence $(\phi_n)\subset X$ 
 satisfying $\mathcal{J}(\phi_n)\to c<\overline{c}$ and $\mathcal{J}'(\phi_n)\to 0$ 
 in $X'$ has a convergent subsequence.
\item  For each $n\in \mathbb{N}$, there exists a set $A_n\subset \Gamma_n$ such that 
$\sup_{\phi\in A_n}\mathcal{J}(\phi)<0$.
\end{enumerate}
Then, there exists a sequence $(\phi_n)$ such that $\mathcal{J}'(\phi_n)=0$, 
$\mathcal{J}(\phi_n)<0$, and $\phi_n\to 0$ in $X$.
\end{lemma}

\section{Auxiliary problem}
The classic variational theory cannot be used to the energy functional 
$\mathcal{J}_\lambda$ since it is not Fr\'echet differentiable due to the singular term. 
So, to prove our main result, we introduce the perturbed problem
\begin{equation}\label{weaksystem}
\begin{gathered}
m(\mathcal{N}(\phi,\psi))\mathcal{L}_\mathcal{K}^p(\phi)=\lambda a(x)(\phi^+
 +\varepsilon)^{-\gamma_1}+\frac{\theta_1}{p^\star_{N,s}}(\phi^+)^{\theta_1-1}
 (\psi^+)^{\theta_2} \quad \text{in } \mathcal{D} \\
m(\mathcal{N}(\phi,\psi))\mathcal{L}_\mathcal{K}^p(\psi)=\lambda b(x)(\psi^+
 +\varepsilon)^{-\gamma_2}+\frac{\theta_2}{p^\star_{N,s}}(\phi^+)^{\theta_1}
 (\psi^+)^{\theta_2-1} \quad \text{in } \mathcal{D} \\
\phi=\psi=0 \quad \text{in } \mathbb{R}^N\setminus \mathcal{D},
\end{gathered}
\end{equation}
where $\varepsilon\in (0,1)$ is a real number. The energy functional associated with 
problem \eqref{weaksystem} is 
$$
\mathcal{J}_{\lambda,\varepsilon}(\phi,\psi)
=\frac{1}{p}\mathscr{M}(\mathcal{N}(\phi,\psi))
-\lambda \Phi_\varepsilon(\phi,\psi)-\Psi(\phi,\psi),
$$
where
$$
\Phi_\varepsilon(\phi,\psi)=\frac{1}{1-\gamma_1}\int_\mathcal{D} a(x)\Big[(\phi^+
+\varepsilon)^{1-\gamma_1}-\varepsilon^{1-\gamma_1}\Big]\,dx
+\frac{1}{1-\gamma_2}\int_\mathcal{D}b(x)\Big[(\psi^+
+\varepsilon)^{1-\gamma_2}-\varepsilon^{1-\gamma_2}\Big]\,dx,
$$
It is easy to see that the functional $\mathcal{J}_{\varepsilon,\lambda}$ is of 
$C^1$ in $X$, and for any $(\phi,\psi)\in X$ and $(\varphi_1,\varphi_2)\in X$, we have
\begin{equation} \label{p11}
\begin{aligned}
&\mathcal{J}'_{\lambda,\varepsilon}(\phi,\psi)\cdot (\varphi_1,\varphi_2) \\
&=m(\mathcal{N}(\phi,\psi))\int_{\mathbb{R}^{2N}}|\phi(x)-\phi(y)|^{p-2}(\phi(x)-\phi(y))
 (\varphi_1(x)-\varphi_1(y))\mathcal{K}(x,y)\,dx\,dy \\
&\quad +m(\mathcal{N}(\phi,\psi))\int_{\mathbb{R}^{2N}}|\psi(x)-\psi(y)|^{p-2}(\psi(x)-\psi(y))
 (\varphi_2(x)-\varphi_2(y))\mathcal{K}(x,y)\,dx\,dy \\
&\quad -\lambda \int_\mathcal{D} a(x)(\phi^++\varepsilon)^{-\gamma_1}\varphi_1\,dx
 -\lambda\int_\mathcal{D}b(x)(\psi^++\varepsilon)^{-\gamma_2}\varphi_2\,dx  \\
&\quad -\frac{\theta_1}{p^\star_{N,s}}\int_\mathcal{D}(\phi^+)^{\theta_1-1}
(\psi^+)^{\theta_2}\varphi_1\,dx-\frac{\theta_2}{p^\star_{N,s}}
\int_\mathcal{D}(\phi^+)^{\theta_1}(\psi^+)^{\theta_2-1}\varphi_2\,dx.
\end{aligned}
\end{equation}
Now, we present the following lemma concerning the Palais Smale condition, which will
play a crucial role in the proof of our main result.

\begin{definition} \rm
A functional $\mathcal{J}$ is said to be satisfying the Palais-Smale condition at 
level $L_\lambda\in \mathbb{R}$ if any sequence $(\phi_n,\psi_n)_{n\in\mathbb{N}}$ in 
$X$ with
\[
\mathcal{J}(\phi_n,\psi_n)\to L_\lambda\quad \text{and}\quad
\mathcal{J}'(\phi_n,\psi_n)\to 0\quad \text{in }  X'\quad \text{as } n\to +\infty,
\]
possesses a convergent subsequence in $X$.
\end{definition}

\begin{lemma}\label{lemma1} 
Under assumptions {\rm (A1)--(A3)}, there exists $\lambda_0>0$ such that for any 
$\lambda\in (0,\lambda_0)$, the functional $\mathcal{J}_{\varepsilon,\lambda}$ 
satisfies the Palais-Smale condition at any level $L_\lambda<0$.
\end{lemma}

\begin{proof}
Let $(\phi_n,\psi_n)$ be a sequence of $X$ such that
\begin{equation}\label{PS}
\mathcal{J}_{\lambda,\varepsilon}(\phi_n,\psi_n)\to L_\lambda\quad\text{and}\quad
\mathcal{J}'_{\lambda,\varepsilon}(\phi_n,\psi_n)\to 0\quad\text{in } X'\quad\text{as } 
n\to +\infty,
\end{equation}
Due to the degenerate type of system \eqref{weaksystem}, two situations have to be 
considered: either
$$
\inf_{n\in \mathbb{N}} \mathcal{N}(\phi_n,\psi_n)=\delta>0,\quad
 \text{and}\quad
 \inf_{n\in \mathbb{N}} \mathcal{N}(\phi_n,\psi_n)=0.
$$
For this, we divide the proof into two cases.
\smallskip

\noindent\textbf{Case 1:} $\inf_{n\in \mathbb{N}} \mathcal{N}(\phi_n,\psi_n)=\delta>0$. 
At the beginning, we show that $(\phi_n,\psi_n)$ is bounded in $X$. By assumption 
(A1), there exists $\kappa_\delta>0$ such that
\begin{equation}\label{k}
m(\mathcal{N}(\phi_n,\psi_n))\geq \kappa_\delta\quad \text{for all } n\in\mathbb{N}.
\end{equation}
By Equation \eqref{PS}, and \eqref{k} with (A2), and the elementary inequality
\begin{equation}\label{elementary}
(\alpha+\beta)^{1-\gamma}-\beta^{1-\gamma}\leq \alpha^{1-\gamma},\quad
 \alpha,\beta\geq 0,\; \gamma\in (0,1),
\end{equation}
we have
\begin{equation} \label{K000}
\begin{aligned}
L_\lambda+o_n(1)
&=\mathcal{J}_{\lambda,\varepsilon}(\phi_n,\psi_n)-\frac{1}{p^\star_{N,s}}
 \mathcal{J}'_{\lambda,\varepsilon}(\phi_n,\psi_n)\cdot (\phi_n,\psi_n) \\
&\quad =\frac{1}{p}\mathscr{M}(\mathcal{N}(\phi_n,\psi_n))-\frac{\lambda}{1-\gamma_1}
 \int_\mathcal{D} a(x)\Big[(\phi_n^++\varepsilon)^{1-\gamma_1}
 -\varepsilon^{1-\gamma_1}\Big]\,dx \\
&\quad -\frac{\lambda}{1-\gamma_2}\int_\mathcal{D} b(x)\Big[(\psi_n^+
 +\varepsilon)^{1-\gamma_2}-\varepsilon^{1-\gamma_2}\Big]\,dx
 -\frac{1}{p^\star_{N,s}}\int_\mathcal{D}  (\phi_n^+)^{\theta_1}(\psi_n^+)^{\theta_2}\,dx
 \\
&\quad -\frac{1}{p^\star_{N,s}}m(\mathcal{N}(\phi_n,\psi_n))\mathcal{N}(\phi_n,\psi_n)
 +\frac{\lambda}{p^\star_{N,s}}\int_\mathcal{D} a(x)(\phi_n^++\varepsilon)^{-\gamma_1}
 \phi_n\,dx  \\
&\quad +\frac{\lambda}{p^\star_{N,s}}\int_\mathcal{D}b(x)(\psi_n^+
 +\varepsilon)^{-\gamma_2}\psi_n\,dx+\frac{1}{p^\star_{N,s}}\int_\mathcal{D}
 (\phi_n^+)^{\theta_1}(\psi_n^+)^{\theta_2}\,dx \\
&\geq(\frac{\kappa_\delta}{p\sigma}-\frac{1}{p^\star_{N,s}})\| \phi_n,\psi_n\|^p
 -\lambda(\frac{1}{p^\star_{N,s}}+\frac{1}{1-\gamma_1})
  \int_\mathcal{D} a(x)(\phi_n^+)^{1-\gamma_1}\,dx\\
&\quad -\lambda(\frac{1}{p^\star_{N,s}}+\frac{1}{1-\gamma_2})
 \int_\mathcal{D} b(x)(\psi_n^+)^{1-\gamma_2}\,dx
\end{aligned}
\end{equation}
On the other hand, by using  H\"older's inequality, and \eqref{sobolevinequaity}
  we obtain
\begin{equation} \label{K00}
\begin{aligned}
\int_\mathcal{D} a(x)(\phi_n^+)^{1-\gamma_1} dx
&\leq |a|_{L^\infty(\mathcal{D})}|\mathcal{D}|^\frac{p+\gamma_1-1}{p}
 |\phi_n|_{L^p(\mathcal{D})}^{1-\gamma_1} \\
&\leq |a|_{L^\infty(\mathcal{D})}|\mathcal{D}|^\frac{p+\gamma_1-1}{p}
 C_p^{1-\gamma_1}|\phi_n|_{E_0}^{1-\gamma_1}\\
&\leq C_a \| \phi_n,\psi_n\|^{1-\gamma_1},
\end{aligned}
\end{equation}
and
\begin{equation}  \label{K01}
\begin{aligned}
\int_\mathcal{D} b(x)(\psi_n^+)^{1-\gamma_2} dx
&\leq |b|_{L^\infty(\mathcal{D})}|\mathcal{D}|^\frac{p+\gamma_2-1}{p}
 |\psi_n|_{L^p(\mathcal{D})}^{1-\gamma_2} \\
&\leq |b|_{L^\infty(\mathcal{D})}|\mathcal{D}|^\frac{p+\gamma_2-1}{p}
 C_p^{1-\gamma_2}|\psi_n|_{E_0}^{1-\gamma_2}\\
&\leq C_b \| \phi_n,\psi_n\|^{1-\gamma_2},
\end{aligned}
\end{equation}
where
$$
C_a=|a|_{L^\infty(\mathcal{D})}|\mathcal{D}|^\frac{p+\gamma_1-1}{p}C_p^{1-\gamma_1}\quad
\text{and}\quad
C_b=|b|_{L^\infty(\mathcal{D})}|\mathcal{D}|^\frac{p+\gamma_2-1}{p}C_p^{1-\gamma_2}.
$$
Now, combining equations \eqref{K00} and \eqref{K01} with \eqref{K000}, we obtain
the boundedness of $(\phi_n,\psi_n)$. Now, since $X$ is a reflexive space, there exist
$(\phi,\psi)\in X$ such that, up to a subsequence (still denoted by $(\phi_n,\psi_n)$),
$(\phi_n,\psi_n)\rightharpoonup (\phi,\psi)$ weakly in $X$, strongly in
$ L^\alpha(\mathcal{D})\times L^\beta(\mathcal{D})$, for any $(\alpha,\beta)\in [1,p^\star_{N,s})^2$,
almost every where in $\mathcal{D}\times\mathcal{D}$ as $n\to +\infty$. Furthermore,
there exist two positive real numbers $\eta$ and $\mu$ such that
\begin{equation}\label{K0}
\mathcal{N}(\phi_n,\psi_n)\to \mu,\quad \text{and}\quad
\int_\mathcal{D} |\phi^+_n-\phi^+|^{\theta_1}|\psi^+_n-\psi^+|^{\theta_2}\,dx\to \eta,
\end{equation}
as $n\to +\infty$. Clearly, $\mu>0$ since we are in the case where $\delta>0$.
Therefore, the weak convergence of $(\phi_n,\psi_n))$ give that the sequence
$(\Phi_n,\Psi_n)$ defined in
$\mathbb{R}^{2N}\setminus \text{diag}\lbrace \mathbb{R}^{2N}\rbrace$ by
\begin{gather*}
A_n(x,y)=|\phi_n(x)-\phi_n(y)|^{p-2}(\phi_n(x)-\phi_n(y))\mathcal{K}(x,y)^\frac{1}{p'},\\
B_n(x,y)=|\psi_n(x)-\psi_n(y)|^{p-2}(\psi_n(x)-\psi_n(y))\mathcal{K}(x,y)^\frac{1}{p'},
\end{gather*}
is bounded in $L^{p'}(\mathbb{R}^{2N})$, as well as
\begin{gather*}
A_n(x,y)\longrightarrow |\phi(x)-\phi(y)|^{p-2}(\phi(x)
-\phi(y))\mathcal{K}(x,y)^\frac{1}{p'} \text{a.e. in}\ \mathbb{R}^{2N}, \\
B_n(x,y)\longrightarrow |\psi(x)-\psi(y)|^{p-2}(\psi(x)
 -\psi(y))\mathcal{K}(x,y)^\frac{1}{p'} \text{a.e. in}\ \mathbb{R}^{2N}.
\end{gather*}
Here $p'=\frac{p}{p-1}$ which is defined the conjugate of $p$.
Then, going if necessary to a further subsequence, we have
$$
A_n(x,y)\longrightarrow |\phi(x)-\phi(y)|^{p-2}(\phi(x)-\phi(y))
\mathcal{K}(x,y)^\frac{1}{p'}
$$
and
$$
B_n(x,y)\longrightarrow |\psi(x)-\psi(y)|^{p-2}(\psi(x)-\psi(y))
\mathcal{K}(x,y)^\frac{1}{p'}
$$
in $L^{p'}(\mathbb{R}^{2N})$. Consequently, as $n\to +\infty$, we obtain
\begin{equation} \label{K1}
\begin{aligned}
&\int_{\mathbb{R}^{2N}}|\phi_n(x)-\phi_n(y)|^{p-2}(\phi_n(x)
 -\phi_n(y))(\phi_1(x)-\phi(y))\mathcal{K}(x,y)\,dx\,dy\\
&= \int_{\mathbb{R}^{2N}}|\phi(x)-\phi(y)|^{p-2}(\phi(x)-\phi(y))(\phi(x)
 -\phi(y))\mathcal{K}(x,y)\,dx\,dy+o_n(1)
\end{aligned}
\end{equation}
and
\begin{equation} \label{K2}
\begin{aligned}
&\int_{\mathbb{R}^{2N}}|\psi_n(x)-\psi_n(y)|^{p-2}(\psi_n(x)
 -\psi_n(y))(\psi(x)-\psi(y))\mathcal{K}(x,y)\,dx\,dy\\
&= \int_{\mathbb{R}^{2N}}|\psi(x)-\psi(y)|^{p-2}(\psi(x)
 -\psi(y))(\psi(x)-\psi(y))\mathcal{K}(x,y)\,dx\,dy+o_n(1).
\end{aligned}
\end{equation}
On the other hand, by H\"older's inequality and Equation \eqref{S}, we have
\begin{align*}
\int_\mathcal{D}\Big|( \phi^+_n)^{\theta_1-1} (\psi^+_n)^{\theta_2}
 \Big|^\frac{p^\star_{N,s}}{p^\star_{N,s}-1}\,dx
&\leq \Big(\int_\mathcal{D}|\phi_n|^{{p^\star_{N,s}}}\,dx
 \Big)^\frac{(\theta_1-1)}{p^\star_{N,s}-1}\Big(\int_\mathcal{D}
 |\psi_n|^{{p^\star_{N,s}}}\,dx\Big)^\frac{\theta_2 }{p^\star_{N,s}-1}\\
&\leq \Big(S_\star^{-\frac{(\theta_1-1)}{p^\star_{N,s}-1}}\|
 \phi_n\|_{E_0}^\frac{(\theta_1-1)}{p^\star_{N,s}}\Big)
 \Big(S_\star^{-\frac{\theta_2p^\star_{N,s}}{p^\star_{N,s}-1}}
 \| \psi_n\|_{E_0)}^\frac{\theta_2p^\star_{N,s}}{p^\star_{N,s}}\Big)
\leq  C,
\end{align*}
for some $C>0$. Similarly, one has
\[
\int_\mathcal{D}\Big|(\phi^+_n)^{\theta_1}( \psi^+_n)^{\theta_2-1}
\Big|^\frac{p^\star_{N,s}}{p^\star_{N,s}-1}\,dx\leq C.
\]
Therefore, as $n\to +\infty$, we have
\begin{gather}\label{K3}
\int_\mathcal{D} ( \phi^+_n)^{\theta_1-1} (\psi^+_n)^{\theta_2}\phi(x)\,dx
=\int_\mathcal{D} ( \phi^+)^{\theta_1} (\psi^+)^{\theta_2}\,dx+o_n(1), \\
\label{K4}
\int_\mathcal{D} ( \phi^+_n)^{\theta_1} (\psi^+_n)^{\theta_2-1}\psi(x)\,dx
=\int_\mathcal{D} ( \phi^+)^{\theta_1} (\psi^+)^{\theta_2}\,dx+o_n(1).
\end{gather}
In addition, we have
\begin{gather}\label{K5}
\Big|\int_\mathcal{D}a(x)(\phi_n^++\varepsilon)^{-\gamma_1}(\phi_n-\phi)\,dx \Big|
\leq |a|_{L^\infty(\mathcal{D})}\varepsilon^{-\gamma_1} \int_\mathcal{D}|\phi_n
 -\phi|\,dx  =o_n(1), \\
\label{K6}
\Big|\int_\mathcal{D}b(x)(\psi_n^++\varepsilon)^{-\gamma_2}(\psi_n-\psi)\,dx \Big|
\leq |b|_{L^\infty(\mathcal{D})}\varepsilon^{-\gamma_2} \int_\mathcal{D}|\psi_n
 -\psi|\,dx =o_n(1),
\end{gather}
as $n\to +\infty$.
Combining equations \eqref{PS} with \eqref{K0}-\eqref{K6}, we obtain
\begin{align*}
o_n(1)
&=\mathcal{J}'_{\lambda,\varepsilon}(\phi_n,\psi_n)\cdot(\phi_n-\phi,\psi_n-\psi)\\
&=m(\mathcal{N}(\phi_n,\psi_n))\Big(\mathcal{N}(\phi_n,\psi_n)-\mathcal{N}
 (\phi,\psi)\Big)\\
&\quad -\int_\mathcal{D}a(x)(\phi_n^++\varepsilon)^{-\gamma_1}(\phi_n-\phi)\,dx\\
&\quad -\int_\mathcal{D}b(x)(\psi_n^++\varepsilon)^{-\gamma_2}(\psi_n-\psi)\,dx\\
&\quad -\int_\mathcal{D}(\phi_n^+)^{\theta_1}(\psi_n^+)^{\theta_2}\,dx
 +\int_\mathcal{D}(\phi^+)^{\theta_1}(\psi^+)^{\theta_2}\,dx+ o_n(1)  \\
&=m(\mu)\| \phi_n-\phi,\psi_n-\psi\|^p
 -\int_\Omega |\phi^+_n(x)-\phi^+(x)|^{\theta_1}|\psi^+_n(x)
 -\psi^+(x)|^{\theta_2}\,dx+o_n(1),
\end{align*}
thanks to Lemmas \ref{l1} and \ref{l3}. Consequently,
\begin{equation}\label{K11}
m(\mu)\| \phi_n-\phi,\psi_n-\psi\|^p=\int_\Omega |\phi^+_n(x)
-\phi^+(x)|^{\theta_1}|\psi^+_n(x)-\psi^+(x)|^{\theta_2}\,dx+o_n(1).
\end{equation}
When $\eta=0$, since $m$ possesses a unique zero at $0$ and $\mu>0$, Equation \eqref{K11}
gives
$$
\|\phi_n-\phi,\psi_n-\psi\|^p=o_n(1),
$$
concluding the proof.  Assuming, by contradiction that $\eta>0$.
From equation \eqref{K11}, it follows that
\begin{equation}\label{K22}
\eta=m(\mu)\left(\mu-\|\phi,\psi\|^p\right).
\end{equation}
In addition, by equation \eqref{S} and Equation \eqref{K11}, we have:
\begin{equation}\label{K33}
\eta^\frac{p^\star_{N,s}-p}{p^\star_{N,s}}\geq m(\mu)S^\frac{p}{p^\star_{N,s}}.
\end{equation}
Now, utilizing  \eqref{K22} and \eqref{K33}, we have
\begin{equation}\label{K44}
m(\mu)S^\frac{p}{p^\star_{N,s}}\leq \eta^\frac{p^\star_s-p}{p^\star_s}=\left[m(\mu)\left(\mu-\| u,v\|^p\right)\right]^\frac{p^\star_{N,s}-p}{p^\star_{N,s}}.
\end{equation}
The latter gives
\begin{equation}\label{K54}
\left(\mu-\| \phi,\psi \|^p\right)^\frac{p^\star_{N,s}-p}{p}\geq m(\mu) S.
\end{equation}
Since the exact behavior of $m$ is unknown, we need to consider two additional cases,
either $\mu\in (0,1)$ or $\mu\geq 1$.  To proceed, we divide the proof of the first
 case into two subcases.
\smallskip

\noindent\textbf{Subcase $\mu\in (0,1)$:} By condition (A3) and inequality
\eqref{K54}, we have
\[
\mu^\frac{p^\star_{N,s}-p}{p} \geq \big(\mu-\| u,v\|^p\big)^\frac{p^\star_{N,s}-p}{p}\\
\geq  m(\mu) S
\geq  m_0 \mu^{\sigma-1}S.
\]
Since $\theta \sigma<p_{N,s}^\star$, it follows that:
\begin{equation}\label{K55}
\mu\geq \big(m_0 S\big)^\frac{p}{p_{N,s}^\star-\sigma p}.
\end{equation}
By utilizing condition (A3), inequalities \eqref{K33}, and \eqref{K55}, one has
\begin{equation} \label{K66}
\begin{aligned}
\eta\geq \Big( m(\mu)S^\frac{p}{p^\star_{N,s}}
\Big)^{\frac{p^\star_s}{p^\star_{N,s}-p}}
&\geq \Big(m_0\mu^{\sigma-1}S^\frac{p}{p^\star_{N,s}} \Big)^{\frac{p^\star_{N,s}}{p^\star_{N,s}-p}}, \\
&\geq \Big[m_0\big(S m_0\big)^\frac{(\sigma-1)p}{p^\star_{N,s}-\sigma p}S^\frac{p}{p^\star_{N,s}} \Big]^{\frac{p^\star_{N,s}}{p^\star_{N,s}-p}},\\
&= (m_0^{p^\star_s}S^{\sigma p})^\frac{1}{p^\star_{N,s}-\sigma p}.
\end{aligned}
\end{equation}
By considering assumptions (A1) and (A3), we obtain the  inequality
\begin{equation}\label{K77}
\mathscr{M}(\mu)\geq \frac{1}{\sigma}m(\mu)\mu\geq \frac{1}{\sigma}m_0\mu^\sigma.
\end{equation}
Let $\sigma_0>\sigma$ be such that $p\sigma_0<p^\star_{N,s}$.
From equations \eqref{PS}, \eqref{elementary}, \eqref{K00}, \eqref{K01}, \eqref{K77},
we deduce that
\begin{equation} \label{eq01}
\begin{aligned}
L_\lambda+o_n(1)
&=\mathcal{J}_{\lambda,\varepsilon}(\phi_n,\psi_n)-\frac{1}{p\sigma_0}\mathcal{J}_{\lambda,\varepsilon}'(\phi_n,\psi_n)\cdot(\phi_n,\psi_n) \\
&=\frac{1}{p}\mathscr{M}(\mathcal{N}(\phi_n,\psi_n))-\frac{\lambda}{1-\gamma_1}\int_\mathcal{D} a(x)\Big[(\phi_n^++\varepsilon)^{1-\gamma_1}-\varepsilon^{1-\gamma_1}\Big]\,dx\\
&\quad-\frac{\lambda}{1-\gamma_2}\int_\mathcal{D} b(x)\Big[(\psi_n^++\varepsilon)^{1-\gamma_2}-\varepsilon^{1-\gamma_2}\Big]\,dx-\frac{1}{p^\star_{N,s}}\int_\mathcal{D}  (\phi_n^+)^{\theta_1}(\psi_n^+)^{\theta_2}\,dx \\
&\quad-\frac{1}{p\sigma_0}m(\mathcal{N}(\phi_n,\psi_n))\mathcal{N}(\phi_n,\psi_n) \\
&\quad+\frac{\lambda}{p\sigma_0}\int_\mathcal{D} a(x)(\phi_n^++\varepsilon)^{-\gamma_1}\phi_n\,dx+\frac{\lambda}{p\sigma_0}\int_\mathcal{D}b(x)(\psi_n^++\varepsilon)^{-\gamma_2}\psi_n\,dx \\
&\quad+\frac{1}{p\sigma_0}\int_\mathcal{D}(\phi_n^+)^{\theta_1}(\psi_n^+)^{\theta_2}\,dx \\
&\geq\frac{1}{p}\mathscr{M}(\mu)-\frac{1}{p\sigma_0}m(\mu)\mu+\Big(\frac{1}{p\sigma_0}-\frac{1}{p^\star_{N,s}}\Big)\int_\mathcal{D}(\phi_n^+)^{\theta_1}(\psi_n^+)^{\theta_2}\,dx \\
&\quad-2\lambda|a|_{L^\infty(\mathcal{D})} \int_\mathcal{D} (\phi_n^+)^{1-\gamma_1}\,dx-2\lambda|b|_{L^\infty(\mathcal{D})} \int_\mathcal{D} (\psi_n^+)^{1-\gamma_2}\,dx+o_n(1), \\
&\geq\frac{1}{p}\mathscr{M}(\mu)-\frac{1}{p\sigma_0}m(\mu)\mu+\Big(\frac{1}{p\sigma_0}-\frac{1}{p^\star_{N,s}}\Big)\int_\mathcal{D}(\phi_n^+)^{\theta_1}(\psi_n^+)^{\theta_2}\,dx \\
&\quad-2\lambda C_a \| \phi_n,\psi_n\|^{1-\gamma_1}-2\lambda C_b \| \phi_n,\psi_n\|^{1-\gamma_2}+o_n(1) \\
&\geq\Big(\frac{1}{\sigma}-\frac{1}{\sigma_0}\Big)\frac{m_0}{p}\mu^\theta+\Big(\frac{1}{p\sigma_0}-\frac{1}{p^\star_{N,s}}\Big)\eta-\lambda C \max\lbrace \mu^\frac{1-\gamma_1}{p},\mu^\frac{1-\gamma_2}{p} \rbrace+o_n(1).
\end{aligned}
\end{equation}
Here $C=2(C_a+C_b)$. Without loos of generality, we assume may that
$\gamma_1\leq \gamma_2$. Thus, since $\mu\in (0,1)$, one has
\[
L_\lambda\geq\Big(\frac{1}{\sigma}-\frac{1}{\sigma_0}\Big)
\frac{m_0}{p}\mu^\theta+\Big(\frac{1}{p\sigma_0}
-\frac{1}{p^\star_{N,s}}\Big)\eta-\lambda C \mu^\frac{1-\gamma_2}{p} .
\]
By employing the Young inequality, we obtain
\begin{align*}
L_\lambda
&\geq \Big(\frac{1}{\theta}-\frac{1}{\sigma_0}\Big)
 \frac{m_0}{p}\mu^\sigma+\Big(\frac{1}{p\sigma_0}-\frac{1}{p^\star_{N,s}}\Big)\eta-\lambda C\mu^\frac{1-\gamma_2}{p},\\
&\geq \Big(\frac{1}{\sigma}-\frac{1}{\sigma_0}\Big)
 \frac{m_0}{p}\mu^\sigma+\Big(\frac{1}{p\sigma_0}-\frac{1}{p^\star_{N,s}}\Big)\eta-\Big(\frac{1}{\sigma}-\frac{1}{\sigma_0}\Big)\frac{m_0}{p}\mu^\sigma\\
&\quad -( \lambda C)^\frac{\sigma p}{\theta p-(1-\gamma_2)}
 \Big((\frac{1}{\sigma}-\frac{1}{\sigma_0})\frac{m_0}{p}\Big)^{-\frac{1-\gamma_2}{\theta p-(2-\gamma_1-\gamma_2)}} ,\\
&\geq\Big(\frac{1}{p\sigma_0}-\frac{1}{p^\star_{N,s}}\Big)\eta
 -( \lambda C)^\frac{\sigma p}{\sigma p-(1-\gamma_2)}
 \Big((\frac{1}{\sigma}-\frac{1}{\sigma_0})
 \frac{m_0}{p}\Big)^{-\frac{2-\gamma_1-\gamma_2}{\sigma p-(1-\gamma_2)}} ,
\end{align*}
Let
\begin{equation}\label{lambda1}
\lambda'=\frac{1}{C}\Big((\frac{1}{\sigma}
 -\frac{1}{\sigma_0})\frac{m_0}{p}\Big)^{\frac{1-\gamma_2}{\sigma p}}
 \Big[(\frac{1}{\sigma p}-\frac{1}{p^\star_{N,s}})(m_0^{p^\star_s}
 S^{p^\star_{N,s}\sigma})^\frac{1}{p^\star_{N,s}
 -\sigma p}\Big]^\frac{\sigma p-(1-\gamma_2)}{\sigma p}.
\end{equation}
Then, we deduce  that for any $\lambda<\lambda'$,
\begin{equation*}
0>L_\lambda\geq\Big(\frac{1}{\sigma_0p}
-\frac{1}{p^\star_{N,s}}\Big)\eta-(\lambda C)^\frac{\sigma p}{\sigma p-(1-\gamma_2)}
((\frac{1}{\sigma}-\frac{1}{\sigma_0})\frac{m_0}{p})^{-\frac{1-\gamma_2}
{\sigma p-(1-\gamma_2)}} >0.
\end{equation*}
Thanks to equation \eqref{K66}, we obtain our contradiction concluding the proof
of the first sub-case.
\smallskip

\noindent\textbf{Subcase $\mu\geq 1$:} 
From condition (A1) for $\delta=1$, there exists $\kappa_\delta>0$ such that
\begin{equation}\label{K78}
m(\mu)\geq \kappa_\delta \quad \text{for all } \mu\geq \delta.
\end{equation}
Combining equation \eqref{K78} with equation \eqref{K33}, we obtain
\begin{equation}\label{K88}
\eta\geq \Big(m(\mu)S^\frac{p}{p^\star_{N,s}}\Big)^\frac{p^\star_{N,s}}{p^\star_{N,s}-p}
 \geq \Big(\kappa_\delta S^\frac{p}{p^\star_{N,s}}\Big)^\frac{p^\star_s}
 {p^\star_{N,s}-p}.
\end{equation}
Proceeding as earlier, with equations \eqref{K78} and \eqref{K88} we obtain
\[
L_\lambda \geq \Big(\frac{1}{\sigma}-\frac{1}{\sigma_0}\Big)\frac{\kappa_\delta}{p}\mu
 +\Big(\frac{1}{p\sigma_0}-\frac{1}{p^\star_{N,s}}\Big)\eta
  -\lambda C \mu^\frac{1-\gamma_2}{p},\\
\geq \Big(\frac{1}{p\sigma_0}-\frac{1}{p^\star_{N,s}}\Big)\eta
 -\Big[\frac{\lambda C}{\big(\frac{1}{\sigma}-\frac{1}{\sigma_0}\big)
 \frac{\kappa_\delta}{p}}\Big]^\frac{p}{p-(1-\gamma_2)}.
\]
Let
\begin{equation}\label{lambda2}
\lambda''=\Big(\frac{1}{\sigma}-\frac{1}{\sigma_0}\Big)
\frac{\kappa_\delta}{pC }\Big[\Big(\frac{1}{p\sigma_0}
-\frac{1}{p^\star_{N,s}}\Big)\Big(\kappa_1S_\star^\frac{p}{p^\star_s}
 \Big)^\frac{p^\star_s}{p^\star_{N,s}-p}\Big]^\frac{p-(1-\gamma_2)}{p}.
\end{equation}
It follows from equation \eqref{K88} that for any $\lambda<\lambda''$, one has
\[
0>L_\lambda\geq\Big(\frac{1}{p\sigma_0}-\frac{1}{p^\star_{N,s}}\Big)
\Big(\kappa_\delta S^\frac{p}{p^\star_{N,s}}\Big)^\frac{p^\star_{N,s}}{p^\star_{N,s}-p}
-\Big[\frac{\lambda C}{\big(\frac{1}{\sigma}-\frac{1}{\sigma_0}\big)
\frac{\kappa_\delta}{p}}\Big]^\frac{p}{p-(1-\gamma_)}>0.
\]
Hence, we still have a contradiction, which concludes the proof of the first case.
\smallskip

\noindent\textbf{Case 2:}
 $\inf_{n\in \mathbb{N}} \mathcal{N}(\phi_n,\psi_n)=0$. In this case, we have two 
possibilities. Either $(0,0)$ is an accumulation point for the sequence 
$(\phi_n,\psi_n)$, and so that, there exists a subsequence of $(\phi_n,\psi_n)$ 
(still denoted by $(\phi_n,\psi_n)$), that strongly converges to $(\phi,\psi)=(0,0)$, 
or $(0,0)$ is an isolated point of $(\phi_n,\psi_n)$.

The first case cannot occur because it would give that the trivial solution $(0,0)$ 
is a critical point at the level $L_\lambda$. However, this is impossible since we have 
$\mathcal{J}_{\lambda,\varepsilon}(0,0)=L_\lambda<0$. Thus, only the latter case 
can occur, which means that there exists a subsequence, (still denoted by 
$(\phi_n,\psi_n)$), such that $\inf_{n\in\mathbb{N}}\mathcal{N}(\phi_n,\psi_n)=\delta>0$. 
We can proceed as before by considering this subsequence.

This completes the proof of the second case. 
Hence, $\mathcal{J}_{\lambda,\varepsilon}$ satisfies the Palais-Smale condition at any 
level $L_\lambda<0$ for any $\lambda\in (0,\lambda_0)$, where 
$\lambda_0=\min \lbrace\lambda',\lambda'' \rbrace$.
\end{proof}

\section{Infinitely many small solutions with negative energy}

Let us note that the functional $\mathcal{J}_{\lambda,\varepsilon}$ is not bounded 
from below in $X$. Indeed, let $(\phi,\psi)\in X$ with $\|(\phi,\psi)\|\leq 1$. 
From (A1), we have $m(t)>0$ for any $t>0$. On the other hand, condition (A2) gives 
that $\frac{M(t)}{\mathscr{M}(t)}\leq\frac{\sigma}{t}$. 
Integrating over $[1,t]$ with $t>1$, we obtain
\[
\mathscr{M}(1)t^\sigma\geq \mathscr{M}(t), \quad \text{for all } t\geq 1.
\]
The latter gives
\begin{align*}
\mathcal{J}_{\lambda,\varepsilon}(t\phi,t\psi)&=\frac{1}{p}\mathscr{M}(\mathcal{N}(t\phi,t\psi))-\frac{\lambda}{1-\gamma_1}\int_\mathcal{D} a(x)\Big[(t\phi^++\varepsilon)^{1-\gamma_1}-\varepsilon^{1-\gamma_1}\Big]\,dx\\
&\quad -\frac{\lambda}{1-\gamma_2}\int_\mathcal{D} b(x)\Big[(t\psi^+
 +\varepsilon)^{1-\gamma_2}-\varepsilon^{1-\gamma_2}\Big]\,dx
 -\frac{1}{p^\star_{N,s}}\int_\mathcal{D}  (t\phi^+)^{\theta_1}(t\psi^+)^{\theta_2}\,dx\\
&\leq \frac{\mathscr{M}(1)}{p}t^{p\sigma}
 +\frac{\lambda\varepsilon^{1-\gamma_1}}{1-\gamma_1}\int_\mathcal{D}a(x)\,dx
 +\frac{\lambda\varepsilon^{1-\gamma_2}}{1-\gamma_2} \int_\mathcal{D}b(x)\,dx\\
&\quad -\frac{1}{p^\star_{N,s}}t^{\theta_1+\theta_2}
 \int_\mathcal{D}  (\phi^+)^{\theta_1}(\psi^+)^{\theta_2}\,dx,
\end{align*}
so that $\mathcal{J}_{\lambda,\varepsilon}(t\phi,t\psi)\to -\infty$ as $t\to +\infty$, 
thanks to the fact that $p\sigma<\theta_1+\theta_2$. Hence, to obtain the existence 
of weak solutions to problem \eqref{weaksystem}, we introduce a suitable truncated 
functional related to $\mathcal{J}_{\lambda,\varepsilon}$ that satisfies the 
conditions of Lemma \ref{kajikiya}.

By using equations \eqref{S}, \eqref{K00} and \eqref{K01}, for any 
$(\phi,\psi)\in X$, we have
\begin{align*}
\mathcal{J}_{\lambda,\varepsilon}(\phi,\psi)
&= \frac{1}{p}\mathscr{M}(\mathcal{N}(\phi,\psi))
 -\frac{\lambda}{1-\gamma_1}\int_\mathcal{D} a(x)\Big[(\phi^+
 +\varepsilon)^{1-\gamma_1}-\varepsilon^{1-\gamma_1}\Big]\,dx\\
&\quad - \frac{\lambda}{1-\gamma_2}\int_\mathcal{D} b(x)
 \Big[(\psi^++\varepsilon)^{1-\gamma_2}-\varepsilon^{1-\gamma_2}\Big]\,dx\\
&\quad - \frac{1}{p^\star_{N,s}}\int_\mathcal{D}  (\phi^+)^{\theta_1}
 (\psi^+)^{\theta_2}\,dx\\
&\geq \frac{1}{p}\mathscr{M}(\|\phi,\psi\|^p)
 - \lambda c_a\| \phi,\psi\|^{1-\gamma_1}-\lambda c_b\| \phi,\psi\|^{1-\gamma_2}
 -\frac{1}{Sp^\star_{N,s}}\| \phi,\psi\|^{p^\star_{N,s}}.
\end{align*}
On the other hand, by (A1) and (A2), we have:

If $\| \phi,\psi \|\leq 1$, then
$$\mathcal{J}_{\lambda,\varepsilon}(\phi,\psi)\geq \frac{\mathscr{M}(1)}{p}\| \phi,\psi \|^{\sigma p}- \lambda c_a\| \phi,\psi\|^{1-\gamma_1}-\lambda c_b\| \phi,\psi\|^{1-\gamma_2}-\frac{S^{-1}}{p^\star_{N,s}}\| \phi,\psi\|^{p^\star_{N,s}}.$$

If $\| \phi,\psi \|> 1$, then
$$
\mathcal{J}_{\lambda,\varepsilon}(\phi,\psi)\geq \frac{1}{p}\kappa\| \phi,\psi\|^p
- \lambda c_a\| \phi,\psi\|^{1-\gamma_1}-\lambda c_b\| \phi,\psi\|^{1-\gamma_2}
-\frac{S^{-1}}{p^\star_{N,s}}\| \phi,\psi\|^{p^\star_{N,s}}.
$$
We define
\[
H(\xi)=
\begin{cases}
\frac{\mathscr{M}(1)}{p}\xi^{p\theta}-\lambda c_a\xi^{1-\gamma_1}
 -\lambda c_b\xi^{1-\gamma_2}-\frac{S^{-1}}{p^\star_{N,s}}\xi^{p^\star_{N,s}} 
 & \text{if } \xi\leq 1\\[4pt]
\frac{\kappa}{p} \xi^p-\lambda c_a\xi^{1-\gamma_1}-\lambda c_b\xi^{1-\gamma_2}
-\frac{S^{-1}}{p^\star_{N,s}}\xi^{p^\star_{N,s}}& \text{if } \xi> 1.
\end{cases}
\]
It is easy to see that for any $\lambda>0$ sufficiently small enough, there exist 
$\xi_1,\xi_2\in (0,1)$ with $\xi_1<\xi_2$ such that $H(\xi_1)=H(\xi_2)=0$ and
\[
H(\xi) \begin{cases}
<0 & \text{if } 0<\xi<\xi_1\\
>0 & \text{if } \xi_1<\xi<\xi_2\,.
\end{cases}
\]
Now, following the same approach as in \cite{t9}, we introduce the following truncated 
functional $\mathcal{T}_\lambda:X\to \mathbb{R}$ defined as
\begin{align*}
\mathcal{T}_\lambda(\phi,\psi)
&=\frac{1}{p}\mathscr{M}(\mathcal{N}(\phi,\psi))
 -\frac{\lambda}{1-\gamma_1}\int_\mathcal{D} a(x)\Big[(\phi^++\varepsilon)^{1-\gamma_1}
 -\varepsilon^{1-\gamma_1}\Big]\,dx\\
&\quad -\frac{\lambda}{1-\gamma_2}\int_\mathcal{D} b(x)\Big[(\psi^+
 +\varepsilon)^{1-\gamma_2}-\varepsilon^{1-\gamma_2}\Big]\,dx
 -\frac{\chi(\| \phi,\psi\|)}{p^\star_{N,s}}\int_\mathcal{D} 
  (\phi^+)^{\theta_1}(\psi^+)^{\theta_2}\,dx
\end{align*}
Here, $\chi:\mathbb{R}+\to [0,1]$ is a non-increasing smooth function such that
$\chi(\xi)=0$ if $\xi\geq \xi_2$ and $\chi(\xi)=1$ if $\xi\leq \xi_1$.

By the construction of $\mathcal{T}_\lambda$ with Lemma \ref{PS}, it can be easily 
verified that $\mathcal{T}_\lambda$ possesses the following properties:
\begin{lemma}\label{L1}
\begin{enumerate}
\item The functional $\mathcal{T}_\lambda$ is of $C^1$, even, and bounded from below 
on $X$.

\item If $\mathcal{T}_\lambda(\phi,\psi)< 0$, then $\| \phi,\psi\|< \xi_1$, and 
$\mathcal{T}_\lambda(\phi,\psi)=\mathcal{J}_{\lambda,\varepsilon}(\phi,\psi)$.

\item For any $\lambda\in (0,\lambda_0)$, $\mathcal{T}_\lambda$ satisfies a local 
Palais-Smale condition for $L_\lambda<0$.
\end{enumerate}
\end{lemma}

Let us define the set
\begin{equation*}
\Sigma_k=\lbrace A\in X\setminus\lbrace 0\rbrace: A \text{ is closed}, 
A=-A, \gamma(A)\geq k\rbrace,
\end{equation*}
where $\gamma(A)$ denotes the genus of $A$, (see Definition \ref{genus}). 
For $k\in\mathbb{N}$, we define the number
\[
C_k=\inf_{A\in \Sigma_k}\sup_{(\phi,\psi)\in A}\mathcal{T}_\lambda(\phi,\psi).
\]
The significance of the real number $C_k$ is that it provides a lower bound for the 
critical values of $\mathcal{T}_\lambda$ restricted to certain subsets of $X$ with 
higher Lusternik-Schnirelmann category.

\begin{lemma}\label{L2}
Assuming that  {\rm (A1)--(A3)} hold. Then 
for each $\lambda>0$ sufficiently small enough and $k\in \mathbb{N}$, $C_k<0$,
 each $C_k$ is a critical value of $\mathcal{T}_\lambda$.
\end{lemma}

\begin{proof}
Let $\xi\in (0,\xi_1)$ and $(\phi,\psi)\in X$ with $\| \phi,\psi\|=1$, we have 
$\chi(\| \xi\phi,\xi\psi\|)=1$. Now, let's fix $k\in \mathbb{N}$ and define 
$X^{(k)}$ as a $k$-dimensional subspace of 
$\mathcal{C}^\infty_0(\mathcal{D})\times \mathcal{C}^\infty_0(\mathcal{D})$. 
Since all the norms in $X^{(k)}$ are equivalent, there exists $r_k\in (0,1)$ 
such that for any $(\phi,\psi)\in X^{(k)}$, we have
$$
\| \phi,\psi\|\leq r_k \implies |\phi,\psi|_{L^\infty(\mathcal{D}\times\mathcal{D})}
\leq 1.
$$
Let us consider the set
$$
\mathscr{S}_{r_k}^{(k)}=\lbrace (\phi,\psi)\in X^{(k)}: \| \phi,\psi\|=r_k\rbrace.
$$
We choose $(\phi^{(k)},\psi^{(k)})\in \mathscr{S}_{r_k}^{(k)}$ such that 
$\phi^{(k)}> 0$, $\psi^{(k)}> 0$, and $\xi\in (0,\xi_1)$. Then,
\begin{align*}
&\mathcal{T}_{\lambda}(\xi\phi^{(k)},\xi\psi^{(k)}) \\
&=\frac{1}{p}\mathscr{M}(\mathcal{N}(\xi\phi^{(k)},\xi\psi^{(k)}))
 -\frac{\lambda}{1-\gamma_1}\int_\mathcal{D} a(x)
 \Big[(\xi\phi^{(k)}+\varepsilon)^{1-\gamma_1}-\varepsilon^{1-\gamma_1}\Big]\,dx\\
&\quad -\frac{\lambda}{1-\gamma_2}\int_\mathcal{D} b(x)
 \Big[(\xi\psi^{(k)}+\varepsilon)^{1-\gamma_2}-\varepsilon^{1-\gamma_2}\Big]\,dx
 -\frac{\xi^{p^\star_{N,s}}}{p^\star_{N,s}}\int_\mathcal{D} 
  (\phi^{(k)})^{\theta_1}(\psi^{(k)})^{\theta_2}\,dx\\
&\leq\frac{\xi^p}{p}\Big(\max_{z\in(0,\xi_1)}m(z)
 \Big)\| \phi^{(k)},\psi^{(k)}\|^p\\
&\quad -\frac{\lambda}{1-\gamma_1}\int_\mathcal{D} a(x)
 \Big[(\xi\phi^{(k)}+\varepsilon)^{1-\gamma_1}-\varepsilon^{1-\gamma_1}\Big]\,dx\\
&\quad -\frac{\lambda}{1-\gamma_2}\int_\mathcal{D} b(x)
\Big[(\xi\psi^{(k)}+\varepsilon)^{1-\gamma_2}-\varepsilon^{1-\gamma_2}\Big]\,dx\\
&\leq\frac{\xi^p}{p}\Big(\max_{z\in(0,\xi_1)}m(z)
 \Big)\| \phi^{(k)},\psi^{(k)}\|^p\\
&\quad -\xi^\frac{1-\gamma_1}{p}\varepsilon^\frac{(1-\gamma_1)(p-1)}{p}
 \int_\mathcal{D}a(x)\big(\phi^{(k)}\big)^\frac{1-\gamma_1}{p}\,dx
 -\xi^\frac{1-\gamma_2}{p}\varepsilon^\frac{(1-\gamma_2)(p-1)}{p}
 \int_\mathcal{D}b(x)\big(\psi^{(k)}\big)^\frac{1-\gamma_2}{p}\,dx,
\end{align*}
thanks to the elementary inequality
$$
(z+t)^{1-\gamma}-t^{1-\gamma}
\geq (1-\gamma)z^\frac{1-\gamma}{p}t^\frac{(1-\gamma)(p-1)}{p},\quad
 p>1,\; t>0,\; z>0 \text{ large enough}.
$$
Since $p>\frac{1-\gamma_1}{p}$ and $p>\frac{1-\gamma_2}{p}$, then we can find 
$\xi_k\in (0,t_1)$ and $\tau_k$ such that
\[
\mathcal{T}_{\lambda}(\xi_k\phi^{(k)},\xi_k\psi^{(k)})\leq -\tau_k<0,\quad 
\forall (\phi^{(k)},\psi^{(k)})\in \mathscr{S}_{r_k}^{(k)}, 
\]
which gives
\[
\mathcal{T}_{\lambda}(\phi^{(k)},\psi^{(k)})\leq -\tau_k<0,\quad 
\forall (\phi^{(k)},\psi^{(k)})\in \mathscr{S}_{\xi_k\rho_k}^{(k)}.
\]
Thus, $C_k<0$ for all $k\in\mathbb{N}$. By applying Lemma \ref{L2}, we can deduce that 
$\mathcal{I}_\lambda$ is bounded from below and satisfies the Palais-Smale condition 
at the level $C_k$ for any $\lambda>0$ small enough. Furthermore, 
Proposition \ref{prop} states that $\gamma(\mathscr{S}_{r_k}^{(k)})=k$. 
Consequently, according to Lemma \ref{kajikiya}, we can conclude that each $C_k$, 
$k\in \mathbb{N}$, is a critical value of $\mathcal{T}_{\lambda}$. 
Consequently, according to $(2)$ of Lemma \ref{L1}, $\mathcal{J}_{\lambda,\varepsilon}$ 
admits a sequence of critical points 
$(\phi_{k,\varepsilon},\psi_{k,\varepsilon})\subset X$ that converges to zero. 
Now, to show the positivity of the solutions 
$(\phi_{k,\varepsilon},\psi_{k,\varepsilon})$, we replace the test function 
$(\varphi_1,\varphi_2)$ in the equation \eqref{p11} by 
$\phi_{k,\varepsilon}^-=\max\lbrace -\phi_{k,\varepsilon},0\rbrace$, 
$\psi_{k,\varepsilon}^-=\max\lbrace -\psi_{k,\varepsilon},0\rbrace$  and using the 
elementary inequality
$$
(a-b)(a^--b^-)\leq -(a^--b^-)^2, 
$$
we obtain $\| \phi_{k,\varepsilon}^-,\psi_{k,\varepsilon}^-\|=0$ implying that 
$(\phi_{k,\varepsilon})$ and $(\psi_{k,\varepsilon})$ are two nonnegative functions. 
By applying the maximum principle (Proposition $2.17$, \cite{max}), we conclude that 
$(\phi_{k,\varepsilon},\psi_{k,\varepsilon})$ is a sequence of positive solutions for 
system \eqref{weaksystem}. This completes the proof.
\end{proof}

\section{Proof of our main result}

In this section, we will show that system \eqref{System} possesses a sequence of 
nontrivial weak solutions $(\phi_k,\psi_k)$ in the space $X$ as a limit of the 
solutions of problem \eqref{weaksystem} obtained in the previous section. 
Let $\lambda\in (0,\lambda_0)$ be small enough and for $k\in\mathbb{N}$, 
let $(\phi_{k,\varepsilon},\psi_{k,\varepsilon})_{\varepsilon>0}$ be a family of
 positive weak solutions of problem \eqref{weaksystem}.
\smallskip

\noindent\textbf{Case 1:}
$\inf_{\varepsilon>0} \mathcal{N}(\phi_{k,\varepsilon},\psi_{k,\varepsilon})=\delta_k>0$.
From  (A1), there exists $\kappa_k>0$ such that
\[
m(\mathcal{N}(\phi_{k,\varepsilon},\psi_{k,\varepsilon}))\geq \kappa_k,\quad
 \text{for all } \varepsilon>0.
\]
From this, (A2), and equations \eqref{K00} and \eqref{K01}, we have
\begin{align*}
C_k+o_\varepsilon(1)
&=\mathcal{J}_{\lambda,\varepsilon}(\phi_{k,\varepsilon},\psi_{k,\varepsilon})-\frac{1}{p^\star_{N,s}}\mathcal{J}'_{\lambda,\varepsilon}(\phi_{k,\varepsilon},\psi_{k,\varepsilon})\cdot (\phi_{k,\varepsilon},\psi_{k,\varepsilon})\\
&=\frac{1}{p}\mathscr{M}(\mathcal{N}(\phi_{k,\varepsilon},\psi_{k,\varepsilon}))-\frac{\lambda}{1-\gamma_1}\int_\mathcal{D} a(x)\Big[(\phi_{k,\varepsilon}+\varepsilon)^{1-\gamma_1}-\varepsilon^{1-\gamma_1}\Big]\,dx\\
&\quad-\frac{\lambda}{1-\gamma_2}\int_\mathcal{D} b(x)\Big[(\psi_{k,\varepsilon}+\varepsilon)^{1-\gamma_2}-\varepsilon^{1-\gamma_2}\Big]\,dx-\frac{1}{p^\star_{N,s}}\int_\mathcal{D}  \phi_{k,\varepsilon}^{\theta_1}\psi_{k,\varepsilon}^{\theta_2}\,dx\\
&\quad-\frac{1}{p^\star_{N,s}}m(\mathcal{N}(\phi_{k,\varepsilon},\psi_{k,\varepsilon}))\mathcal{N}(\phi_{k,\varepsilon},\psi_{k,\varepsilon})+\frac{\lambda}{p^\star_{N,s}}\int_\mathcal{D} a(x) (\phi_{k,\varepsilon}+\varepsilon)^{-\gamma_1}\phi_{k,\varepsilon}\,dx \\
&\quad+\frac{\lambda}{p^\star_{N,s}}\int_\mathcal{D}b(x)(\psi_{k,\varepsilon}+\varepsilon)^{-\gamma_2}\psi_{k,\varepsilon}\,dx+\frac{1}{p^\star_{N,s}}\int_\mathcal{D}\phi_{k,\varepsilon}^{\theta_1}\psi_{k,\varepsilon}^{\theta_2}\,dx\\
&\geq(\frac{\kappa_\delta}{p\sigma}-\frac{1}{p^\star_{N,s}})\| \phi_{k,\varepsilon},\psi_{k,\varepsilon}\|^p-\lambda(\frac{1}{p^\star_{N,s}}+\frac{1}{1-\gamma_1})\int_\mathcal{D} a(x)(\phi_n^+)^{1-\gamma_1}\,dx\\
&\quad-\lambda(\frac{1}{p^\star_{N,s}}+\frac{1}{1-\gamma_2})\int_\mathcal{D} b(x)(\psi_n^+)^{1-\gamma_2}\,dx\\
&\geq(\frac{\kappa_\delta}{p\sigma}-\frac{1}{p^\star_{N,s}})\| \phi_{k,\varepsilon},\psi_{k,\varepsilon}\|^p-\lambda(\frac{1}{p^\star_{N,s}}+\frac{1}{1-\gamma_1})C_a \| \phi_{k,\varepsilon},\psi_{k,\varepsilon}\|^{1-\gamma_1}\\
&\quad-\lambda(\frac{1}{p^\star_{N,s}}+\frac{1}{1-\gamma_2})C_b \| 
\phi_{k,\varepsilon},\psi_{k,\varepsilon}\|^{1-\gamma_2}
\end{align*}
This gives the boundedness of 
$(\phi_{k,\varepsilon},\psi_{k,\varepsilon})_{\varepsilon>0}$. 
Since the space $X$ is reflexive, up to a subsequence, still denoted by 
$(\phi_{k,\varepsilon},\psi_{k,\varepsilon})_{\varepsilon>0}$, there exists 
$(\phi_k,\psi_k)\in X$ such that, 
$(\phi_{k,\varepsilon},\psi_{k,\varepsilon})\rightharpoonup (\phi_k,\psi_k)$ 
weakly in $X$, $(\phi_{k,\varepsilon},\psi_{k,\varepsilon})\to (\phi_k,\psi_k)$ 
strongly in 
$L^\alpha(\mathcal{D})\times L^\beta(\mathcal{D})$  for all $(\alpha,\beta)\in [1,p^\star_{N,s})^2$,
$(\phi_{k,\varepsilon},\psi_{k,\varepsilon})\to (\phi_k,\psi_k)$  a.e.\
in $\mathcal{D}\times \mathcal{D}$, as $\varepsilon\to 0^+$. In addition, there exist
$\mu,\eta>0$ and $f_1,f_2\in L^1(\mathcal{D})$ such that $\phi_{k,\varepsilon}\leq f_1$, 
$\psi_{k,\varepsilon}\leq f_2$, and
\begin{equation}\label{eq22}
\mathcal{N}(\phi_{k,\varepsilon},\psi_{k,\varepsilon})\to \mu,\quad \text{and}\quad
\int_\mathcal{D} |\phi_{k,\varepsilon}-\phi_k|^{\theta_1}|\psi_{k,\varepsilon}
-\psi_k|^{\theta_2}\,dx\to \eta.
\end{equation}
We aim to show that 
$(\phi_{k,\varepsilon},\psi_{k,\varepsilon})\rightharpoonup (\phi_k,\psi_k)$ 
strongly in $X$. We observe that
\begin{gather*}
|a(x) (\phi_{k,\varepsilon}+\varepsilon)^{-\gamma_1}\phi_{k,\varepsilon}|
\leq a(x)\phi_{k,\varepsilon}^{1-\gamma_1}, \\
|b(x) (\psi_{k,\varepsilon}+\varepsilon)^{-\gamma_2}\psi_{k,\varepsilon}|
\leq b(x)\psi_{k,\varepsilon}^{1-\gamma_2}, 
\end{gather*}
a.e.\ in $\mathcal{D}$, so by the Vitali convergence theorem, we obtain
\begin{gather}\label{eq33}
\int_\mathcal{D}a(x) (\phi_{k,\varepsilon}+\varepsilon)^{-\gamma_1}
\phi_{k,\varepsilon}\,dx=\int_\mathcal{D}a(x)\phi_k^{1-\gamma_1}\,dx,\\
\label{eq44}
\int_\mathcal{D}b(x) (\psi_{k,\varepsilon}+\varepsilon)^{-\gamma_2}
\psi_{k,\varepsilon}\,dx=\int_\mathcal{D}b(x)\psi_k^{1-\gamma_2}\,dx.
\end{gather}
On the other hand, a simple calculation in \eqref{weaksystem}, we have
\begin{gather*}
m(\mathcal{N}(\phi_{k,\varepsilon},\psi_{k,\varepsilon}))
\mathcal{L}_p(\phi_{k,\varepsilon})=\lambda(\phi_{k,\varepsilon}
 +\varepsilon)^{-\gamma_1}+\frac{\theta_1}{p^\star_{N,s}}\phi_{k,\varepsilon}^{\theta_1-1}\psi_{k,\varepsilon}^{\theta_2}\geq \min\lbrace \frac{\lambda}{2^{\gamma_1}} ,\frac{\theta_1}{p^\star_{N,s}}\rbrace  \\
m(\mathcal{N}(\phi_{k,\varepsilon},\psi_{k,\varepsilon}))\mathcal{L}_p
 (\psi_{k,\varepsilon})=\lambda(\psi_{k,\varepsilon}+\varepsilon)^{-\gamma_2}
 +\frac{\theta_2}{p^\star_{N,s}}\phi_{k,\varepsilon}^{\theta_1}
 \psi_{k,\varepsilon}^{\theta_2-1}\geq \min\lbrace  
 \frac{\lambda}{2^{\gamma_2}},\frac{\theta_1}{p^\star_{N,s}}\rbrace.
\end{gather*}
Therefore, since 
$\inf_{\varepsilon>0} \mathcal{N}(\phi_{k,\varepsilon},\psi_{k,\varepsilon})=\delta_k>0$ 
and using the strong maximum principle (see \cite{max}), there exist 
$\mathcal{D}_0\subset \mathcal{D}$, $\mathcal{D}_1\subset \mathcal{D}$ and a constants 
$d_0>0$, $d_1$ that are independent of $\varepsilon$ such that for any $\varepsilon>0$, 
we have
\begin{equation}\label{eq66}
\phi_{k,\varepsilon}\geq d_0 \text{ a.e. in }  \mathcal{D}_0, \quad \text{and}\quad 
\psi_{k,\varepsilon}\geq d_1 \text{ a.e.in } \mathcal{D}_1.
\end{equation}
Now, let us consider a function $\varphi_1\in \mathcal{C}^\infty_0(\mathcal{D}$ with 
$\operatorname{supp}(\varphi_1)\subset \mathcal{D}_0$ and a function 
$\varphi_2\in \mathcal{C}^\infty_0(\mathcal{D}$ with 
$\operatorname{supp}(\varphi_2)\subset \mathcal{D}_0$.
 We observe that from \eqref{eq66} we have
\begin{gather*}
|a(x)(\phi_{k,\varepsilon}+\varepsilon)^{-\gamma_1}\varphi_1|
\leq a(x) d_0^{-\gamma_1}|\varphi_1|,\quad \text{a.e. in } \mathcal{D}_0,\\
|b(x)(\psi_{k,\varepsilon}+\varepsilon)^{-\gamma_2}\varphi_2|
\leq b(x) d_1^{-\gamma_2}|\varphi_2|,\quad \text{a.e. in } \mathcal{D}_1.
\end{gather*}
Consequently, the dominated convergence theorem implies that
\begin{gather}\label{eq77}
\int_\mathcal{D}a(x) (\phi_{k,\varepsilon}+\varepsilon)^{-\gamma_1}\varphi_1\,dx
=\int_\mathcal{D}a(x) \phi_{k}^{-\gamma_1}\varphi_1\,dx+o_\varepsilon(1), \\
\label{eq88}
\int_\mathcal{D}b(x) (\psi_{k,\varepsilon}+\varepsilon)^{-\gamma_2}\varphi_2\,dx
=\int_\mathcal{D}b(x) \psi_{k}^{-\gamma_2}\varphi_2\,dx+o_\varepsilon(1),
\end{gather}
as $\varepsilon\to0^+$. Since $\partial\mathcal{D}$ is continuous, the space 
$\mathcal{C}^\infty_0(\mathcal{D})$ is dense in $E_0(\mathcal{D})$ 
(see \cite[Theorem 6]{density}). 
Therefore, by a standard density argument, equations \eqref{eq77} and \eqref{eq88}
hold true for any $(\varphi_1,\varphi_2)\in X$. Thus, by combining 
\eqref{eq33}-\eqref{eq44} with \eqref{eq77}-\eqref{eq88} with $\varphi_1=\phi_k$ and 
$\varphi_2=\psi_k$, as $\varepsilon\to0^+$, we obtain
\begin{gather*}
\int_\mathcal{D}a(x) (\phi_{k,\varepsilon}+\varepsilon)^{-\gamma_1}
(\phi_{k,\varepsilon}-\phi_k)\,dx=o_\varepsilon(1), \\
\int_\mathcal{D}b(x) (\psi_{k,\varepsilon}+\varepsilon)^{-\gamma_2}(\psi_{k,\varepsilon}
-\psi_k)\,dx=o_\varepsilon(1).
\end{gather*}
Consequently, from \eqref{eq22}, as $\varepsilon\to 0^+$ we have
\begin{align*}
o_\varepsilon(1)
&=\mathcal{J}'_{\varepsilon,\lambda}(\phi_{k,\varepsilon},\psi_{k,\varepsilon})\cdot(\phi_{k,\varepsilon}-\phi_k,\psi_{k,\varepsilon}-\psi_k)\\
&=m(\mathcal{N}(\phi_{k,\varepsilon},\psi_{k,\varepsilon}))\Big(\mathcal{N}(\phi_{k,\varepsilon},\psi_{k,\varepsilon})-\mathcal{N}(\phi_k,\psi_k)\Big)\\
&\quad -\int_\mathcal{D}a(x) (\phi_{k,\varepsilon}+\varepsilon)^{-\gamma_1}(\phi_{k,\varepsilon}-\phi_k)\,dx-\int_\mathcal{D}b(x) (\psi_{k,\varepsilon}+\varepsilon)^{-\gamma_2}(\psi_{k,\varepsilon}-\psi_k)\,dx\\
&\quad -\int_\mathcal{D}(\phi_{k,\varepsilon})^{\theta_1}(\psi_{k,\varepsilon})^{\theta_2}\,dx+\int_\mathcal{D}(\phi_k)^{\theta_1}(\psi_k)^{\theta_2}\,dx+ o_\varepsilon(1)  \\
&=m(\mu)\| \phi_{k,\varepsilon}-\phi_k,\phi_{k,\varepsilon}-\psi_k\|^p-\int_\Omega |\phi_{k,\varepsilon}(x)-\phi_k(x)|^{\theta_1}|\psi_{k,\varepsilon}(x)-\psi_k(x)|^{\theta_2}\,dx+o_\varepsilon(1),
\end{align*}
so that 
$$
m(\mu)\| \phi_{k,\varepsilon}-\phi_k,\phi_{k,\varepsilon}-\psi_k\|^p
=\eta+o_\varepsilon(1).
$$
Now, let us show that $\eta=0$. By contradiction, i.e., we assume that $\eta>0$. 
Arguing as in Lemma \ref{lemma1}, we can show that
\begin{equation}
\eta^\frac{p^\star_{N,s}-p}{p^\star_{N,s}}\geq m(\mu)S^\frac{p}{p^\star_{N,s}}.
\end{equation}
Now, consider $\sigma_0>\sigma$ be such that $p\sigma_0<p_{N,s}^\star$. 
Then, similarly to \eqref{eq01}, as $\varepsilon\to 0^+$ we have
\begin{align*}
C_k+o_\varepsilon(1)&=\mathcal{J}_{\lambda,\varepsilon}(\phi_{k,\varepsilon},\psi_{k,\varepsilon})-\frac{1}{\sigma_0}\mathcal{J}'_{\lambda,\varepsilon}(\phi_{k,\varepsilon},\psi_{k,\varepsilon})\cdot (\phi_{k,\varepsilon},\psi_{k,\varepsilon})\\
&=\frac{1}{p}\mathscr{M}(\mathcal{N}(\phi_{k,\varepsilon},\psi_{k,\varepsilon}))-\frac{\lambda}{1-\gamma_1}\int_\mathcal{D} a(x)\Big[(\phi_{k,\varepsilon}+\varepsilon)^{1-\gamma_1}-\varepsilon^{1-\gamma_1}\Big]\,dx\\
&\quad-\frac{\lambda}{1-\gamma_2}\int_\mathcal{D} b(x)\Big[(\psi_{k,\varepsilon}+\varepsilon)^{1-\gamma_2}-\varepsilon^{1-\gamma_2}\Big]\,dx-\frac{1}{p^\star_{N,s}}\int_\mathcal{D}  \phi_{k,\varepsilon}^{\theta_1}\psi_{k,\varepsilon}^{\theta_2}\,dx\\
&\quad-\frac{1}{\sigma_0}m(\mathcal{N}(\phi_{k,\varepsilon},\psi_{k,\varepsilon}))\mathcal{N}(\phi_{k,\varepsilon},\psi_{k,\varepsilon})+\frac{\lambda}{\sigma_0}\int_\mathcal{D} a(x) (\phi_{k,\varepsilon}+\varepsilon)^{-\gamma_1}\phi_{k,\varepsilon}\,dx \\
&\quad+\frac{\lambda}{p^\star_{N,s}}\int_\mathcal{D}b(x)(\psi_{k,\varepsilon}+\varepsilon)^{-\gamma_2}\psi_{k,\varepsilon}\,dx+\frac{1}{\sigma_0}\int_\mathcal{D}\phi_{k,\varepsilon}^{\theta_1}\psi_{k,\varepsilon}^{\theta_2}\,dx\\
&\geq\frac{1}{p}\mathscr{M}(\mu)-\frac{1}{p\sigma_0}m(\mu)\mu+\Big(\frac{1}{p\sigma_0}-\frac{1}{p^\star_{N,s}}\Big)\int_\mathcal{D} |\phi_{k,\varepsilon}-\phi_k|^{\theta_1}|\psi_{k,\varepsilon}-\psi_k|^{\theta_2}\,dx\\
&\quad+\Big(\frac{1}{p\sigma_0}-\frac{1}{p^\star_{N,s}}\Big)\int_\mathcal{D} |\phi_k|^{\theta_1}|\psi_k|^{\theta_2}\,dx-\lambda C\Big(\| \phi_n,\psi_n\|^{1-\gamma_1}+ \| \phi_n,\psi_n\|^{1-\gamma_2}\Big)+o_\varepsilon(1)\\
&\geq\Big(\frac{1}{\sigma}-\frac{1}{\sigma_0}\Big)\frac{m_0}{p}\mu^\theta+\Big(\frac{1}{p\sigma_0}-\frac{1}{p^\star_{N,s}}\Big)\eta-\lambda C\Big(\| \phi_n,\psi_n\|^{1-\gamma_1}+ \| \phi_n,\psi_n\|^{1-\gamma_2}\Big)+o_\varepsilon(1).\\
\end{align*}
That is,
$$
C_k\geq  \Big(\frac{1}{\sigma}-\frac{1}{\sigma_0}\Big)
\frac{m_0}{p}\mu^\theta+\Big(\frac{1}{p\sigma_0}-\frac{1}{p^\star_{N,s}}\Big)
\eta-\lambda C \max\lbrace \mu^\frac{1-\gamma_1}{p},\mu^\frac{1-\gamma_2}{p} \rbrace.
$$
Without loss of generality, we may assume that $\sigma_1\leq\sigma_2$. 
In particular, for any $\lambda\in (0,\lambda_0)$, we obtain
$$
0>C_k\geq \Big(\frac{1}{\sigma}-\frac{1}{\sigma_0}\Big)\frac{m_0}{p}\mu^\theta
+\Big(\frac{1}{p\sigma_0}-\frac{1}{p^\star_{N,s}}\Big)\eta
-\lambda C \mu^\frac{1-\gamma_2}{p}>0, 
$$
thanks to \eqref{lambda1} and \eqref{lambda2}. Which leads to a contradiction concluding 
$\eta=0$. Hence $(\phi_{k,\varepsilon},\psi_{k,\varepsilon})
\rightharpoonup (\phi_k,\psi_k)$ strongly in $X$.
\smallskip

\noindent\textbf{Case 2:} 
$\inf_{\varepsilon>0} \mathcal{N}(\phi_{k,\varepsilon},\psi_{k,\varepsilon})=0$. 
In this case, we can apply a similar approach to that used in the proof of 
Lemma \ref{L2}, allowing us to deduce that 
$(\phi_{k,\varepsilon},\psi_{k,\varepsilon})\rightharpoonup (\phi_k,\psi_k)$ 
strongly in $X$. This complete the proof of Theorem \ref{Theorem}.

\begin{thebibliography}{10}

\bibitem{l1} Applebaum, D.; 
\emph{L\'evy processes-from probability of finance quantum groups}, 
Not. Am. Math. Soc., 51 (2004), 1336-1347.

\bibitem{Alv} Alves, C. O.;  Figueiredo, G. M.;
\emph{Multi-bump solutions for a Kirchhoff-type problem}, 
Adv. Nonlinear Anal., 2016 5(1). 1-26.

\bibitem{ambr} Ambrosio, V.;
\emph{Concentration phenomena for a class of fractional Kirchhoff equations 
in $\mathbb{R}^N$ with general nonlinearities}, Nonlinear Analysis, 195(2020), 111-761.

\bibitem{Au} Autuori, G.;  Fiscella, A.;  Pucci, P.;
\emph{Stationary  Kirchhoff  problems  involving  a  fractional  elliptic operator 
and a critical nonlinearity}, Nonlinear Anal., 125 (2015). 699-714.

\bibitem{t1} Barrios, B.; Colorado, E.; De Pablo, A.; Sanchez, U.;
\emph{On some critical problems for the fractional Laplacian operator}, 
J. Differ. Equ., 252(2012), 6133-6162 .

\bibitem{t2} Br\"andle, C.; Colorado, E.; de Pablo, A.; S\'anchez, U.;
\emph{A concave-convex elliptic problem involving the fractional Laplacian}, 
Proc. R. Soc. Edinb. A, 143(2013), 39-71 .

\bibitem{bb}  Br\'ezis, H., Lieb, E.;
\emph{A relation between pointwise convergence of functions and convergence of 
functional}, Proc. Am. Math. Soc., 88(1983), 486-490.

\bibitem{caff} Caffarelli, L.;
\emph{Nonlocal equations, drifts and games. In: Nonlinear Partial Differential Equations}.
 Abel Symposia, 7(2012), 37-52.

\bibitem{t3} Chen, J. W.; Deng, B. S.;
\emph{The Nehari manifold for non-local elliptic operators involving concave-convex 
nonlinearities}, Z. Angew. Math. Phys., 66(2015), 1387-1400.

\bibitem{t4} Chen, J. W.; Deng, B. S.;
\emph{The Nehari manifold for a fractional $p$-Laplacian system involving concave-convex 
nonlinearities. Nonlinear Anal.}, Real World Appl., 27(2016), 80-92 .

\bibitem{Col} Colasuonno, F.; Pucci, P.;
\emph{Multiplicity of solutions for $p(x)$-polyharmonic elliptic Kirchhoff equations}, 
Nonlinear Anal.,  74 (2011), 5962-74.

\bibitem{D} Daouas, A.; Louchaich, M.;
\emph{On fractional $p$-Laplacian type equations with general nonlinearities,} 
Turk. J. Math., (2021) 45, 2477-2491.

\bibitem{daoua} Daouas, A.; Louchaich, M.; Saoudi, K.;
\emph{Existence of infinitely many solutions for fractional p-Kirchhoff systems 
involving critical Sobolev nonlinearities}. J. Elliptic Parabol Equ.,
 (2025. DOI 10.1007/s41808-025-00362-3

\bibitem{nez} Di Nezza, E.; Palatucci, G.; Valdinoci, E.;
\emph{Hitchhiker's guide to the fractional Sobolev spaces. Bulletin des Sciences 
Math\'ematiques}, (2012), 136 (5), 521-573.

\bibitem{t5} Dong, X. Y., Bai, Z. B., Zhang, S. Q.;
\emph{Positive solutions to boundary value problems of p-Laplacian with fractional 
derivative}, Bound. Value Probl., (2017),  DOI 10.1186/s13661-016-0735-z

\bibitem{f2}  Figueiredo, G. M.;
\emph{Existence of a positive solution for a Kirchhoff problem type with critical 
growth via truncation argument}, J. Math. Anal. Appl., 401 (2013), 706-713.

\bibitem{z} Fiscella, A.;
\emph{Saddle  point  solutions  for  non-local  elliptic  operators},  
Topol. Methods Nonlinear Anal., 44 (2014), 527-538.

\bibitem{fisc}  Fiscella, A.;  Valdinoci, E.;
\emph{A critical Kirchhoff type problem involving a non local operator}, 
Nonliear Anal, 94 (2014), 156-170.

\bibitem{density} Fiscella, A.; Servadei, R.; Valdinoci, E.;
\emph{Density properties for fractional Sobolev spaces}, 
Ann. Acad. Sci. Fenn. Math., 40 (2015), no. 1, 235–253.

\bibitem{a} Iannizzotto, A.; Liu, S.; Perera, K.; Squassina, M.;
\emph{Existence results for fractional $p$-Laplacian problems via Morse theory}, 
Adv. Calc. Var., 9(2016), 101-125.

\bibitem{n} Kajikiya, R.;
\emph{A critical-point theorem related to the symmetric mountain-pass lemma and its
 applications to elliptic equations}, J. Funct. Anal., 225(2005), 352-370.

\bibitem{o} Liang, S. H.; Repov\u s, D.; Zhang, B. L.;
\emph{On the fractional Schr\"odinger-Kirchhoff equations with electromagnetic 
fields and critical nonlinearity}, Comput. Math. Appl. 75 (2018), 1778-1794.

\bibitem{Liu} Liu, J.;  Liao, J. F.;  Tang, C. L.;
 \emph{Positive  solutions  for  Kirchhoff-type  equations  with  critical  
 exponent in $\mathbb{R}^N$}, J. Math. Anal. Appl., 429 (2015), 1153-72.

\bibitem{Mok}  Mokhtari, A.; Moussaoui, T.;  O'Regan, D.;
\emph{Multiplicity results for an impulsive boundary value problem of 
$p(t)$-Kirchhoff type via critical point theory}, Opusc. Math, (2016) 36(5), 631-649.

\bibitem{Ny} Nyamoradi, N.;
\emph{Existence of three solutions for Kirchhoff nonlocal operators of elliptic 
type Math}, Commun., 18 (2013), 489-502.

\bibitem{Ny1} Nyamoradi, N.; Chung,  N. T.;
\emph{Existence  of  solutions  to  nonlocal  Kirchhoff  equations  of  elliptic 
type via genus theory}, Electron. J. Differ. Equ., 2024 (2014) No. 86, 1-12.

\bibitem{Ny2} Nyamoradi, N.; Teng, K.;
\emph{Existence of solutions for a Kirchhoff-type-nonlocal operators of elliptic 
type Commun}, Pure Appl. Anal., 14 (2015), 361-7.

\bibitem{o2} Ourraoui, A.;
\emph{On a $p$-Kirchhoff problem involving a critical nonlinearity}, 
C. R. Acad. Sci. Paris, Ser. I, 352 (2014), 295-298.

\bibitem{max} Del Pezzo, L. M.; Quaas, A.;
\emph{A Hopf's lemma and a strong minimum principle for the fractional p-Laplacian}. 
Journal of Differential Equations, 263(1), 765-778.

\bibitem{p} Pucci, P.; Saldi, S.;
\emph{Critical stationary Kirchhoff equations in $\mathbb{R}^N$ involving nonlocal operators},
Rev. Mat. Iberoam, 32 (2016), 1-22.

\bibitem{s} Pucci, P.; Xiang, M. Q.; Zhang, B. L.;
\emph{Existence and multiplicity of entire solutions for fractional $p$-Kirchhoff 
equations}, Adv. Nonlinear Anal., 5 (2016), 27-55.

\bibitem{t7} Pucci, P.; Xiang, Q. M.; Zhang, B. L.;
\emph{Multiple solutions for nonhomogeneous Schr\"odinger-Kirchhoff type equations 
involving the fractional $p$-Laplacian in $\mathbb{R}^N$}. Calc. Var., 54 (2015), 2785-2806 .

\bibitem{qui} Qiu, H.;  Xiang, M.;
\emph{Existence  of  solutions  for  fractional $p$-Laplacian  problems  via  
Leray-Schauder's nonlinear alternative}. Boundary Value Problems, (2016), 1-8.

\bibitem{cc} Rabinowitz, P. H.;
\emph{Minimax Methods in Critical Point Theory with Applications to Differential 
Equations}, CBME Regional Conference Series in Mathematics, 65. 
American Mathematical Society, Providence RI (1986).

\bibitem{serv2} Servadei, R.; Valdinoci, E.;
\emph{Mountain Pass solutions for non-local elliptic operators}. Journal of
 Mathematical Analysis and Applications, 389 (2012), 887-898.

\bibitem {serv4} Servadei, R.; Valdinoci, E.;
\emph{Variational methods for non-local operators of elliptic type}, 
Discrete and Continuous Dynamical Systems, 33(2013), 2105-2137.

\bibitem{t9} Wang, L.; Zhang, B. L.;
\emph{Infinitely many solutions for Schr\"odinger-Kirchhoff type equations involving 
the fractional $p$-Laplacian and critical exponent}, 
Electron. J. Differ. Equ., 2016 (2016) No. 339, 1-18.

\bibitem{t10} Wei, Y.; Su, X.;
\emph{Multiplicity of solutions for non-local elliptic equations driven by the 
fractional Laplacian}, Calc. Var. Partial Differ. Equ., 52 (2015), 95-124.

\end{thebibliography}
\end{document}
